% ============================================================
% PT Article M5 — version française
% Miroir FR de m5_bridge.tex (canonique EN)
% ============================================================

\documentclass[12pt,a4paper]{article}
\IfFileExists{pt-common-fr.sty}{\usepackage{pt-common-fr}}{\usepackage{../shared/pt-common-fr}}
\usepackage[margin=2.5cm]{geometry}
\usepackage{setspace}
\usepackage{tikz-3dplot}
\onehalfspacing

\title{\textbf{La théorie de la persistance~V~:\\
le pont de l'arithmétique vers la physique}}
\author{Yan Senez\\\texttt{yan.senez@protonmail.com}}
\date{Avril 2026}

\begin{document}
\maketitle

\begin{abstract}
Ceci est le cinquième et dernier article présentant les fondations
mathématiques de la théorie de la persistance (PT). En s'appuyant
sur la dynamique du crible et les lois de
conservation~\cite{companion_M1}, sur la géométrie de l'information
et l'holonomie~\cite{companion_M2}, sur la théorie de la convergence
et du point fixe~\cite{companion_M3}, et sur les structures
mathématiques étendues~\cite{companion_M4}, nous construisons le pont
formel allant de l'arithmétique des sauts de premiers aux observables
physiques.

\textbf{Six axiomes, tous prouvés comme théorèmes.}
Les axiomes BA0--BA5 codifient le pont. BA0 (le champ dynamique est
la suite des sauts $\{g_n\}$) est un théorème (T0,
\cite{companion_M3})~; BA1--BA4 sont des conséquences structurelles
du noyau arithmétique~; BA5 (le couplage produit) est dérivé de la
dualité de Pontryagin appliquée à la décomposition par théorème
chinois des restes~--- un théorème d'analyse harmonique, et non une
hypothèse physique.

\textbf{Quatre unicités fermant tous les degrés de liberté.}
Le crible est unique (T6, \cite{companion_M2})~; la divergence
$D_{\rm KL}$ est unique (G1, Shore--Johnson, \cite{companion_M2})~;
la métrique de Fisher est unique (G3, \v{C}encov, \cite{companion_M2})~;
le point fixe $\mu^* = 15$ est unique (T5, \cite{companion_M3}). Ces
quatre unicités laissent zéro choix dans la construction~: il existe
exactement une théorie de jauge U(1) sur $(\mathbb{Z}, +, \times)$, et
son couplage est
$\alpha_{\rm crible} = \prod_{p \in \{3,5,7\}} \sin^2\!\theta_p = 1/136{,}28$
(couplage nu~: $0{,}55\%$~; après corrections de polarisation du vide
d'écho~: $0{,}11\,\mathrm{ppb}$~; après la chaîne complète d'habillage
développée dans l'article compagnon~\cite{companion_physics}~:
$0{,}004\,\mathrm{ppb}$, ce qui correspond à
$1/\alpha = 137{,}035\,999\,083$).

\textbf{Limite continue et reconstruction.}
Les axiomes d'Osterwalder--Schrader sont vérifiés uniformément pour
toutes les matrices de transfert finies $T_m$ (la positivité par
réflexion est un théorème algébrique~: argument de la matrice de
Gram). La limite inductive est prouvée par contraction de Buchstab,
ce qui donne des suites de Cauchy de fonctions de Schwinger~; le
théorème de reconstruction d'OS produit une TQC wightmanienne. Le
produit tensoriel infini incomplet de von Neumann fournit une
construction explicite indépendante.

\textbf{Test d'ablation et falsifiabilité.}
La permutation de branche ($q_{+} \leftrightarrow q_{-}$) dégrade les
43 observables d'un facteur $106\times$~--- forte indication que
l'affectation sommet/arête est structurelle, non accidentelle. Le
tableau de correspondance wightmanien (W1--W6 + OS) montre que tout
axiome standard de TQC est satisfait à la fois aux niveaux fini et
continu.

Zéro paramètre ajusté. Cinquième et dernier article de la série.
\end{abstract}

\tableofcontents
\newpage

\section{Introduction}
\label{sec:intro}

Cet article est le cinquième et dernier volet de la série présentant
les fondations mathématiques de la théorie de la persistance (PT).
Les quatre articles précédents ont établi~:

\begin{itemize}[nosep]
\item \textbf{Article~I}~\cite{companion_M1}~: le crible
  d'Ératosthène comme système dynamique sur les résidus des sauts de
  premiers, le théorème des transitions interdites~(T1), le lemme
  d'unicité~(L0) et le théorème fondamental des sauts~(TFS).
\item \textbf{Article~II}~\cite{companion_M2}~: la géométrie
  canonique de l'information (unicité de Shore--Johnson G1, unicité
  de \v{C}encov G3), la formule d'holonomie~(BA3) produisant des
  angles trigonométriques à partir de pure arithmétique, et le
  théorème d'unicité d'Ératosthène~(T6).
\item \textbf{Article~III}~\cite{companion_M3}~: la Formule maîtresse
  et le théorème de convergence spectrale~(T4) prouvant
  $\alpha_k \to 1/2$, le théorème du point fixe~(T5) identifiant
  $\mustar = 15$ comme l'unique point fixe stable, et le théorème
  du champ dynamique (T0, fermeture de BA0) promouvant le postulat
  du champ dynamique au rang de théorème.
\item \textbf{Article~IV}~\cite{companion_M4}~: l'algèbre du crible
  et les transformations PT (persistance, intercalateur, tenseur,
  holonomique, décohérence), la métrique PT, le cadre catégoriel,
  la mécanique complexe sur le cercle de rayon $s = 1/2$, et le
  cadre des mathématiques prédictives.
\end{itemize}

Le présent article construit le \emph{pont}~: le passage formel
de l'arithmétique des sauts de premiers vers les observables
physiques. Le pont comporte trois couches. Premièrement, les six
axiomes BA0--BA5 sont tous prouvés comme théorèmes de théorie
des nombres et de théorie de l'information~--- aucune hypothèse
physique irréductible n'intervient. Deuxièmement, la structure
produit
$\alpha_{\rm crible} = \prod_{p \in \{3,5,7\}} \sin^2\!\theta_p$
est dérivée de la dualité de Pontryagin appliquée à la
factorisation TCR~--- un théorème d'analyse harmonique.
Troisièmement, l'identification physique est forcée par quatre
unicités structurelles (crible unique T6, divergence unique G1,
métrique unique G3, point fixe unique T5) fermant tous les degrés
de liberté, combinées à la vérification de tous les axiomes
structurels de l'ÉDQ (42/42 tests PASS dans l'article
compagnon~\cite{companion_physics}).

Le passage du crible discret à une théorie quantique des champs
continue est démontré~: la contraction de Buchstab fournit des
suites de Cauchy de fonctions de Schwinger, les axiomes
d'Osterwalder--Schrader sont vérifiés à la limite, et le théorème
de reconstruction d'OS produit une TQC wightmanienne. Le produit
tensoriel infini incomplet de von Neumann (ITPS) fournit une
vérification indépendante.

L'article inclut également quatre annexes fournissant le matériel
de référence complet~: le système de classification épistémique
(annexe~A), le diagramme d'architecture mathématique (annexe~B),
le catalogue des transformations PT (annexe~C), et la preuve
complète de la limite inductive et de la reconstruction d'OS
(annexe~D).

L'article compagnon~\cite{companion_physics} applique les
corrections d'habillage qui ferment l'écart nu de $0{,}55\%$ à
$0{,}004\,\mathrm{ppb}$ de précision, et étend le pont aux 43
observables complètes du modèle standard.

\section{Le pont de l'arithmétique vers la physique}
\label{sec:bridge}

\epigraph{Le miracle de l'adéquation du langage des mathématiques
à la formulation des lois de la physique est un don merveilleux que
nous ne comprenons ni ne méritons.}{Eugene Wigner}

% ---- Status Box (R1) ----------------------------------------
\begin{statusbox}
\textbf{OBJECTIF~:}
  Construire le pont systématique entre le noyau arithmétique
  (\cite{companion_M1}--\cite{companion_M3}) et les observables
  physiques, via les six axiomes BA0--BA5 et les quatre propriétés
  d'unicité P1--P4. Dériver la forme produit de $\alphaEM$ à partir
  de la dualité de
  Pontryagin~\cite{Pontryagin1966,Rudin1962} et l'identifier avec la
  constante de structure fine
  physique~\cite{CODATA2018} via une chaîne de quatre unicités
  (T6, G1, G3, T5) fermant tous les degrés de liberté.

\textbf{ENTRÉES~:}
  \cite{companion_M1}--\cite{companion_M2}, $\mustar = 15$
  (\cite{companion_M3}), $s = 1/2$ (\cite{companion_M1}), angles
  d'holonomie $\sinp{p}$ (\cite{companion_M2}), identité de la règle
  de Born (article compagnon~\cite{companion_physics}).

\textbf{RÉSULTAT PRINCIPAL~:}
  \begin{itemize}[nosep]
  \item P1--P4~: la suite des sauts de premiers est l'unique suite
    (dans la classe testée) satisfaisant les quatre propriétés
    structurelles (théorème~\ref{thm:P1P4}).
  \item BA0--BA4~: cinq axiomes sont des théorèmes de
    \cite{companion_M1}--\cite{companion_M3} (BA0 = T0).
  \item BA5~: $\alphaEM = \prod_{p \in \{3,5,7\}}
    \sin^2\!\theta_{p,\rm stat}$~--- forme produit \textbf{dérivée}
    de la dualité de Pontryagin (théorème~\ref{thm:pontryagin})~;
    identification physique via quatre unicités
    (théorème~\ref{thm:reconstruction}).
  \item Le pont PT complet est intérieurement cohérent, avec tous
    les axiomes prouvés comme théorèmes, et falsifiable.
  \end{itemize}

\textbf{STATUT~:}
  \THMTAG~(P1--P4, BA0--BA5, limite inductive).
  Forme produit BA5 via Pontryagin (THM)~;
  limite inductive via contraction de Buchstab + reconstruction d'OS
  (THM, théorème~\ref{thm:inductive_limit}).
  T4 fermé (annihilation spectrale, \cite{companion_M3}).
  \textbf{Clôture R50 (pont thêta)~:} convergence exponentielle
  $\gamma_p \sim q^{0{,}73p}$ quantifiée~; OS1--OS5 vérifiés à la
  limite~; reconstruction wightmanienne sous DIC (W1,W3~: THM~;
  W2,W4~: DER-PHYS).
  19/19 tests PASS (\texttt{test\_theta\_bridge\_R50\_PT.py}).

\textbf{NOTE SUR L'ITPS~:}
  L'ITPS de von Neumann s'applique à des espaces composantes de
  dimension quelconque, y compris croissante ($p-1$). Trois arguments
  indépendants établissent la limite inductive~:
  (a)~la construction ITPS n'exige pas de dimension fixée
  (Guichardet 1966, Araki--Woods 1968)~;
  (b)~le chemin principal de la preuve (étapes~1--3) utilise la
  reconstruction d'OS directement à partir des limites des fonctions
  de Schwinger, contournant entièrement l'ITPS. L'ITPS (étape~4) sert
  de vérification indépendante~;
  (c)~l'argument du pont thêta quantifie le taux de convergence comme
  $q^{cp}$ avec $c \approx 0{,}73$, via la factorisation TCR des
  fonctions de Schwinger et la préservation par matrice de Gram de
  OS3 (\texttt{test\_theta\_bridge\_R50\_PT.py}, 19/19 PASS).
  La transformation modulaire de Jacobi motive la structure
  exponentielle mais n'entre pas dans la chaîne logique.
\end{statusbox}

% ============================================================
\subsection{Pourquoi les angles d'holonomie du crible sont les constantes de couplage physiques}
\label{sec:ch9_why}

Avant de plonger dans la machinerie technique, énonçons l'argument
logique de façon directe. Le lecteur qui comprend cette section comprend
le pont~; les sections suivantes fournissent les preuves.

\begin{tcolorbox}[colback=thmblue!5, colframe=thmblue!60, title={\textbf{L'argument en cinq étapes}}]
\begin{enumerate}[leftmargin=*, label=\textbf{\arabic*.}]
\item \textbf{Le crible crée automatiquement une théorie de jauge.}
  La récurrence $r_{n+1} = (r_n + g_n) \bmod p$ transporte les résidus
  autour du cercle discret $\mathbb{Z}/p\mathbb{Z}$. Il s'agit d'une
  connexion de jauge~--- non par analogie, mais par définition~: c'est
  un transport parallèle sur un fibré principal au-dessus de la suite
  des sauts (\cite{companion_M1}). Aucune hypothèse physique n'est
  importée~; c'est de l'arithmétique pure.

\item \textbf{Cette théorie de jauge a un couplage unique, forcé par quatre unicités.}
  \begin{itemize}[nosep]
  \item Le \emph{crible} est unique (T6, \cite{companion_M2}~:
    Ératosthène est l'unique procédure constructive sur
    $(\mathbb{N}, \times)$).
  \item La \emph{divergence} $\DKL$ est unique (G1,
    \cite{companion_M2}~: axiomes de Shore--Johnson).
  \item La \emph{métrique} (Fisher) est unique (G3,
    \cite{companion_M2}~: monotonie de \v{C}encov).
  \item Le \emph{point fixe} $\mustar = 15$ est unique (T5,
    \cite{companion_M3}~: balayage d'auto-cohérence).
  \end{itemize}
  Au point fixe, chaque premier actif $p \in \{3, 5, 7\}$ contribue
  un angle d'holonomie $\sinp{p} = \delta_p(2 - \delta_p)$ (identité
  algébrique, \cite{companion_M2}). La décomposition TCR est
  additive~; la dualité de Pontryagin transforme les structures
  additives en produits multiplicatifs (un théorème d'analyse
  harmonique, non une hypothèse physique). Par conséquent, le
  couplage total est
  $\alpha_{\rm crible} = \prod_{p \in \{3,5,7\}} \sinp{p}$. Il n'y a
  pas d'autre produit à former~: les quatre unicités laissent zéro
  choix.

\item \textbf{La théorie de jauge du crible satisfait tous les axiomes structurels de l'ÉDQ.}
  Le programme de Feynman (l'article
  compagnon~\cite{companion_physics}) vérifie que cette théorie
  satisfait les identités de Ward, l'unitarité, la symétrie de
  croisement, la statistique du spin, CPT, la conservation de
  l'hélicité, et la $\beta$-fonction correcte~--- tous prouvés à
  partir du crible, non importés de la physique (42/42 tests PASS,
  cinq théorèmes S1--S5).

\item \textbf{Quatre unicités ferment tous les degrés de liberté~: ce
  couplage \emph{est} $\alphaEM$.}
  Le crible est unique (T6). La divergence est unique (G1). La
  métrique est unique (G3). Le point fixe est unique (T5). Ensemble,
  elles laissent \emph{zéro} choix dans la construction~: il existe
  exactement une théorie de jauge U(1) sur $(\mathbb{Z}, +, \times)$
  satisfaisant tous les axiomes structurels, et son couplage est
  $\prod \sinp{p}$. Le programme de Feynman (article
  compagnon~\cite{companion_physics}) confirme que cette théorie
  satisfait les identités de Ward, l'unitarité, le croisement, la
  statistique du spin, CPT, et la $\beta$-fonction correcte (42/42
  tests, cinq théorèmes S1--S5). Le passage au continu est prouvé via
  contraction de Buchstab + reconstruction d'OS
  (théorème~\ref{thm:inductive_limit})~; la métrique de Fisher EST la
  géométrie continue (R50). Les quatre unicités réduisent donc le
  pont à une seule étape ontologique~: l'identification
  $\alpha_{\rm crible} = \alphaEM$.

\item \textbf{Conclusion~: le pont est une unique identification testable.}
  La chaîne d'unicités laisse exactement un degré de liberté~:
  l'énoncé ontologique selon lequel le couplage physique
  \emph{instancie} le couplage du crible. Cet énoncé est étiqueté
  \BRIDGETAG\ (et non \THMTAG)~: il est testé~--- non prouvé~--- par
  l'accord de 43 observables à $0{,}30\%$ de déviation moyenne
  (remarque~\ref{rem:R11}). La force du pont ne réside pas dans la
  nécessité logique mais dans l'absence de tout paramètre libre~: si
  l'identification est vraie, les 43 valeurs sont \emph{forcées}~;
  si elle échoue, l'accord est une coïncidence de probabilité
  $\sim 10^{-6}$.
\end{enumerate}
\end{tcolorbox}

\medskip\noindent
Une confirmation indépendante provient d'une branche complètement
différente des mathématiques~: la voie de la métrique de Fisher
(section~\ref{sec:ch9_fisher_route}) retrouve la même valeur par
double intégration de l'unique métrique riemannienne, en utilisant
la géométrie riemannienne au lieu de l'analyse harmonique. Deux
dérivations, des outils différents, la même réponse.

\medskip\noindent
Le reste de ce chapitre fournit les preuves formelles.

% ============================================================
\subsection{La nature du pont}
\label{sec:ch9_intuition}

Le terme «~pont~» mérite une définition soigneuse. Un pont en PT
est une \emph{chaîne de dérivation}~: il envoie les objets
arithmétiques (classes de résidus, fréquences de sauts, angles
d'holonomie) sur des objets physiques (constantes de couplage,
angles de mélange, rapports de masses) via un ensemble d'axiomes
explicitement énoncés.

Le pont comporte trois couches~:
\begin{enumerate}[nosep]
\item \textbf{Noyau arithmétique (BA0--BA4)~:} tous prouvés comme
  théorèmes de théorie des nombres et de théorie de l'information
  (\cite{companion_M1}--\cite{companion_M3}~; BA0 = T0).
\item \textbf{Structure produit (BA5)~:} dérivée de la dualité de
  Pontryagin appliquée à la factorisation TCR. C'est un théorème
  d'analyse harmonique, non une hypothèse physique.
\item \textbf{Identification physique~:} le couplage du crible EST
  le couplage électromagnétique. Cela découle de la fermeture de
  tous les degrés de liberté par quatre unicités (T6, G1, G3, T5)
  et de la vérification de tous les axiomes structurels de l'ÉDQ
  (section~\ref{sec:ch9_wightman}).
\end{enumerate}

Le pont ne contient aucune hypothèse physique irréductible (BA0 est
lui-même un théorème, T0, \cite{companion_M3}).
L'étape de la règle de Born H3, précédemment identifiée comme une
hypothèse séparée, est une identité algébrique (article
compagnon~\cite{companion_physics}, identité~12.2)~:
$P(r) = |\psi(r)|^2$ lorsque $\psi = \sqrt{P}$ par définition. La
forme produit découle de la théorie des représentations~;
l'identification découle de la reconstruction structurelle.

\begin{remarque}[Test d'ablation~: dégradation d'un facteur $106\times$]
\label{rem:ablation}
Une vérification critique de robustesse est le \textbf{test
d'ablation}~: permuter l'affectation de branche
($\qplus \leftrightarrow \qminus$) et recalculer les 43
observables. L'erreur moyenne saute de $0{,}30\%$ à $31{,}8\%$~---
un facteur de dégradation de $\mathbf{106\times}$. C'est une forte
indication que l'affectation sommet/arête est \emph{structurelle},
non accidentelle~: un motif numérologique survivrait à la permutation
de branche, alors qu'une chaîne de dérivation avec la \emph{mauvaise}
affectation de branche ne peut pas reproduire le MS.
\end{remarque}

% ============================================================
\subsection{Unicité structurelle~: P1--P4}
\label{sec:ch9_P1P4}

\begin{thmbox}{Unicité P1--P4}
\label{thm:P1P4}\leanProved{PT.Bridge.BridgeAxioms.BA2\_GFT\_identity}
Dans la classe de comparaison \{sauts de premiers, nombres
chanceux, composés, géométrique aléatoire\}, la suite des sauts
de premiers $\{g_n = p_{n+1} - p_n\}$ est l'unique suite
satisfaisant les quatre propriétés structurelles~:
\begin{enumerate}[leftmargin=*, label=\textbf{P\arabic*.}]
\item \textbf{Transitions interdites (T1)~:} auto-transitions
  $T[r][r] = 0$ dans la dynamique des résidus mod-$p$.
\item \textbf{Fibration modulaire (TCR)~:}
  $\mathbb{Z}/M\mathbb{Z} \cong \prod_p \mathbb{Z}/p\mathbb{Z}$
  fournit des directions arithmétiques indépendantes.
\item \textbf{Point fixe auto-cohérent (T5)~:}
  $\mustar = \sum_{\text{actifs}} p$ a une unique solution stable.
\item \textbf{Universalité de Mertens~:} la voie de convergence est
  compatible avec la loi de Mertens en $\ln\ln$ asymptotiquement.
\end{enumerate}
Contre-exemples~: nombres chanceux 0/4, composés 0/4, géométrique
aléatoire 0/4. Seuls les premiers~: 4/4.
\end{thmbox}

Ce résultat d'unicité justifie l'utilisation de la suite des sauts
de premiers comme champ dynamique privilégié~--- non par hypothèse,
mais par élimination.

% ============================================================
\subsection{Les six axiomes}
\label{sec:ch9_axioms}

\begin{table}[htbp]
\centering
\caption{Les six axiomes de la théorie de la persistance. Tous sont
des théorèmes (BA0 $=$ T0, \cite{companion_M3}).}
\label{tab:bridge_axioms}
\begin{tabular}{cllcl}
\toprule
\textbf{Axiome} & \textbf{Énoncé} & \textbf{Type} & \textbf{Source} & \textbf{Art.} \\
\midrule
BA0 & Champ $= \{g_n\}$ (DDL unique) & THM & T0 (\cite{companion_M3}) & III \\
BA1 & Observables $= g_n \bmod m$ & DER & Connexion de jauge & I \\
BA2 & $\log_2 m = \DKL + H$ (TFS) & ID & Identité algébrique & I \\
BA3 & $\sinp{p} = \delta_p(2 - \delta_p)$ & THM & Holonomie depuis T6 (unicité) & II \\
BA4 & Actifs~: $\gamp{p} > s$~; $\mustar = 15$ & THM & T5 & III \\
\textbf{BA5} & $\alpha_{\rm crible} = \prod \sinp{p}$ & \textbf{THM} & Pontryagin + reconstruction OS & V \\
 & $\alpha_{\rm crible} \stackrel{!}{=} \alphaEM$ & \BRIDGETAG & Identification ontologique & V \\
\bottomrule
\end{tabular}
\end{table}

Le tableau~\ref{tab:bridge_axioms} liste les six axiomes.
BA0 est un \emph{théorème} (T0, \cite{companion_M3})~: la suite des
sauts de premiers est uniquement déterminée par U1--U4. Dans les
unités canoniques du crible (SCU, R51), la masse de l'électron
$m_e = s = 1/2$ en SCU s'ensuit~--- la valeur SI $0{,}511$~MeV est
un facteur de conversion dimensionnel (la proximité numérique
$0{,}511 \approx 0{,}500$ est coïncidente~; voir l'article
compagnon~\cite{companion_physics}).
BA0--BA5 sont toutes des conséquences dérivées~: BA0 de T0
(\cite{companion_M3}), BA1--BA4 du noyau arithmétique
(\cite{companion_M1}--\cite{companion_M3}), et BA5 de la dualité de
Pontryagin (cet article) combinée à la reconstruction structurelle
(figures~\ref{fig:CRT_pontryagin} et~\ref{fig:CRT_tetrahedron}).

\begin{figure}[htbp]
\centering
\begin{tikzpicture}[
  node distance=1.0cm and 1.2cm,
  box/.style={rectangle, draw, rounded corners=3pt,
              minimum width=3.0cm, minimum height=0.9cm,
              font=\small, align=center, inner sep=4pt},
  addbox/.style={box, draw=thmblue, fill=thmblue!8},
  mulbox/.style={box, draw=predred, fill=predred!8},
  arr/.style={->, >=stealth, thick},
  lbl/.style={font=\scriptsize, midway, fill=white, inner sep=1pt}
]
  % Top: additive (TCR)
  \node[addbox] (crt) {$\mathbb{Z}/210\mathbb{Z}$\\(primorielle du crible)};
  \node[addbox, below=of crt] (z3) {$\mathbb{Z}/3\mathbb{Z}$};
  \node[addbox, left=1.0cm of z3] (z2) {$\mathbb{Z}/2\mathbb{Z}$};
  \node[addbox, right=1.0cm of z3] (z5) {$\mathbb{Z}/5\mathbb{Z}$};
  \node[addbox, right=1.0cm of z5] (z7) {$\mathbb{Z}/7\mathbb{Z}$};
  % CRT arrows
  \draw[arr, thmblue] (crt) -- node[lbl] {TCR $\cong$} (z3);
  \draw[arr, thmblue] (crt) -- (z2);
  \draw[arr, thmblue] (crt) -- (z5);
  \draw[arr, thmblue] (crt) -- (z7);
  % oplus labels
  \node[font=\small, thmblue] at ($(z2.east)!0.5!(z3.west)$) {$\oplus$};
  \node[font=\small, thmblue] at ($(z3.east)!0.5!(z5.west)$) {$\oplus$};
  \node[font=\small, thmblue] at ($(z5.east)!0.5!(z7.west)$) {$\oplus$};
  % Middle: Pontryagin duality
  \node[font=\small, align=center] at ($(z3)!0.5!(z5) + (0,-1.5)$)
    (pont) {\textbf{Dualité de Pontryagin}\\[2pt]$\oplus \;\longrightarrow\; \times$\\[2pt]
    les caractères sont \emph{multiplicatifs}};
  \draw[arr, thick, condorange] ($(z3.south) + (-1.5, -0.2)$) -- (pont);
  \draw[arr, thick, condorange] (pont) -- ($(z3.south) + (-1.5, -2.8)$);
  % Bottom: multiplicative (coupling product)
  \node[mulbox, below=3.2cm of z2] (s2) {$\sin^2\!\theta_2$\\(trivial)};
  \node[mulbox, below=3.2cm of z3] (s3) {$\sin^2\!\theta_3$};
  \node[mulbox, below=3.2cm of z5] (s5) {$\sin^2\!\theta_5$};
  \node[mulbox, below=3.2cm of z7] (s7) {$\sin^2\!\theta_7$};
  % times labels
  \node[font=\small, predred] at ($(s2.east)!0.5!(s3.west)$) {$\times$};
  \node[font=\small, predred] at ($(s3.east)!0.5!(s5.west)$) {$\times$};
  \node[font=\small, predred] at ($(s5.east)!0.5!(s7.west)$) {$\times$};
  % Result
  \node[mulbox, below=1.0cm of $(s3)!0.5!(s5)$, minimum width=6cm]
    (alpha) {$\alphaEM = \displaystyle\prod_{p \in \{3,5,7\}}
             \sin^2\!\theta_{p}(\qplus)$\enspace\DERTAG};
  \draw[arr, predred] (s3) -- (alpha);
  \draw[arr, predred] (s5) -- (alpha);
  \draw[arr, predred] (s7) -- (alpha);
  \draw[arr, predred, dashed] (s2) -- (alpha);
\end{tikzpicture}
\caption{De la TCR au produit de couplage via la dualité de
  Pontryagin. La décomposition additive TCR (bleu, haut) s'envoie
  sur un produit multiplicatif de caractères d'holonomie (rouge,
  bas). Chaque premier actif contribue $\sin^2\!\theta_p$
  indépendamment~; leur produit donne $\alphaEM$. Le passage de
  $\oplus$ à $\times$ est un théorème d'analyse harmonique, non
  une hypothèse physique.}
\label{fig:CRT_pontryagin}
\end{figure}

\begin{figure}[htbp]
\centering
\begin{tikzpicture}[scale=1.0, >=stealth,
  grpnode/.style={draw, thick, rounded corners=3pt, minimum width=1.6cm,
                  minimum height=0.7cm, font=\small},
  ptnode/.style={circle, draw, thick, minimum size=0.6cm, font=\small}]
  % === Top: Z/210Z (sieve primorial) ===
  \node[grpnode, condorange, fill=condorange!8] (Z210) at (0, 4.5)
    {$\mathbb{Z}/210\mathbb{Z}$};
  % === Middle row: three factor groups ===
  \node[grpnode, thmblue, fill=thmblue!8] (Z3) at (-4, 2) {$\mathbb{Z}/3\mathbb{Z}$};
  \node[grpnode, thmblue, fill=thmblue!8] (Z5) at ( 0, 2) {$\mathbb{Z}/5\mathbb{Z}$};
  \node[grpnode, thmblue, fill=thmblue!8] (Z7) at ( 4, 2) {$\mathbb{Z}/7\mathbb{Z}$};
  % === CRT projection arrows ===
  \draw[->, condorange, thick] (Z210) -- (Z3)
    node[midway, above left, font=\scriptsize] {TCR};
  \draw[->, condorange, thick] (Z210) -- (Z5)
    node[midway, right, font=\scriptsize] {TCR};
  \draw[->, condorange, thick] (Z210) -- (Z7)
    node[midway, above right, font=\scriptsize] {TCR};
  % === Additive structure ===
  \draw[thmblue, thick, <->] (Z3) -- (Z5)
    node[midway, below, font=\small] {$\oplus$};
  \draw[thmblue, thick, <->] (Z5) -- (Z7)
    node[midway, below, font=\small] {$\oplus$};
  % === Pontryagin duality arrow ===
  \draw[->, thick, gray, line width=1.2pt] (1.8, 0.8) -- (1.8, -0.2)
    node[midway, right, font=\scriptsize, gray, align=left]
    {Pontryagin\\$\oplus \to \times$};
  % === Bottom row: holonomy angles (multiplicative dual) ===
  \node[ptnode, predred, fill=predred!8] (s3) at (-4, -1) {$\sinp{3}$};
  \node[ptnode, predred, fill=predred!8] (s5) at ( 0, -1) {$\sinp{5}$};
  \node[ptnode, predred, fill=predred!8] (s7) at ( 4, -1) {$\sinp{7}$};
  % === Multiplicative links ===
  \draw[predred, thick] (s3) -- (s5) node[midway, below, font=\small] {$\times$};
  \draw[predred, thick] (s5) -- (s7) node[midway, below, font=\small] {$\times$};
  % === Product result ===
  \node[draw, thick, predred, rounded corners=3pt, fill=predred!8,
        minimum width=2.8cm, minimum height=0.7cm, font=\small\bfseries]
    (alpha) at (0, -3) {$\alphaEM = \prod\sinp{p}$};
  \draw[->, predred, thick] (s3) -- (alpha);
  \draw[->, predred, thick] (s5) -- (alpha);
  \draw[->, predred, thick] (s7) -- (alpha);
  % === Duality arrows linking rows ===
  \draw[->, gray, dashed] (Z3.south) -- (s3.north);
  \draw[->, gray, dashed] (Z5.south) -- (s5.north);
  \draw[->, gray, dashed] (Z7.south) -- (s7.north);
\end{tikzpicture}
\caption{Décomposition TCR et dualité de Pontryagin.
  \emph{Haut}~: la primorielle du crible $\mathbb{Z}/210\mathbb{Z}$
  se projette via TCR sur trois groupes facteurs (additifs, bleu).
  \emph{Bas}~: la dualité de Pontryagin envoie la structure additive
  sur le dual multiplicatif (rouge)~: chaque $\sinp{p}$ est
  l'évaluation du caractère au premier~$p$~; leur produit donne
  $\alphaEM$.}
\label{fig:CRT_tetrahedron}
\end{figure}

% ============================================================
\subsection{La forme produit BA5~: dérivation par Pontryagin}
\label{sec:ch9_BA5}

\begin{thmbox}{Produit de Pontryagin (BA5)}\mTHM
\label{thm:BA5}\leanProved{PT.Anomaly.BargmannBA5.BA5\_product\_identity}
Sous BA0--BA4, le couplage électromagnétique au point fixe a la
forme produit~:
\begin{equation}
  \alphaEM = \prod_{p \in \{3,5,7\}} \sinp{p}(\qplus).
  \label{eq:BA5}
\end{equation}
Ce produit est évalué sur la branche sommet ($\qplus = 13/15$).
\end{thmbox}

\begin{thmbox}{Dérivation par Pontryagin de la forme produit}
\label{thm:pontryagin}\leanProved{PT.Anomaly.BargmannBA5.BA5\_alpha\_sieve\_value}
La forme produit~\eqref{eq:BA5} est dérivée de la chaîne suivante,
chaque étape étant un théorème ou une identité algébrique~:

\begin{enumerate}[leftmargin=*, label=\textbf{Étape \arabic*.}]
\item \textbf{Factorisation TCR (théorème).}
  La primorielle du crible $M = 2 \cdot 3 \cdot 5 \cdot 7 = 210$
  donne, par le théorème chinois des restes (coprimité deux à deux)~:
  \begin{equation}
    \mathbb{Z}/210\mathbb{Z} \;\cong\; \mathbb{Z}/2\mathbb{Z} \oplus
    \mathbb{Z}/3\mathbb{Z} \oplus \mathbb{Z}/5\mathbb{Z} \oplus
    \mathbb{Z}/7\mathbb{Z}.
    \label{eq:CRT_210}
  \end{equation}
  C'est un isomorphisme de groupes \textbf{additifs}. Statut~:
  \THMTAG\ (théorème chinois des restes, classique).

\item \textbf{Dualité de Pontryagin (théorème d'analyse harmonique).}
  Pour tout groupe abélien fini $G = G_1 \oplus G_2 \oplus G_3$, le
  groupe dual $\hat{G} = \mathrm{Hom}(G, \mathbb{T})$ satisfait~:
  \begin{equation}
    \widehat{G_1 \oplus G_2 \oplus G_3} \cong
    \hat{G}_1 \times \hat{G}_2 \times \hat{G}_3.
    \label{eq:pontryagin_factor}
  \end{equation}
  Tout caractère $\chi$ de la somme directe se factorise~:
  $\chi(g_1, g_2, g_3) = \chi_a(g_1) \cdot \chi_b(g_2) \cdot \chi_c(g_3)$.
  C'est un théorème standard~\cite{Rudin1962}~: \emph{les caractères
  des sommes directes additives sont des produits multiplicatifs}.
  Statut~: \THMTAG.

\item \textbf{Localité du crible (corollaire TCR).}
  Le crible au premier $p$ enlève la classe de résidus
  $0 \pmod{p}$ et agit \emph{uniquement} sur la composante
  $\mathbb{Z}/p\mathbb{Z}$ de la décomposition TCR. Les opérations
  à des premiers différents agissent sur des facteurs différents et
  donc commutent. Statut~: \IDTAG.

\item \textbf{Structure de produit tensoriel (théorie des représentations).}
  Par le théorème standard sur les représentations irréductibles des
  produits directs, toute représentation irréductible de
  $G_1 \times G_2 \times G_3$ est de la forme
  $\rho = \rho_1 \otimes \rho_2 \otimes \rho_3$.
  Combiné à l'étape~3, l'état du crible au point fixe $\mustar$ se
  décompose en état produit~:
  \begin{equation}
    |\Psi\rangle = |\psi_3\rangle \otimes |\psi_5\rangle
    \otimes |\psi_7\rangle.
    \label{eq:product_state}
  \end{equation}
  Ceci est \textbf{dérivé} de la structure de groupe, non supposé.
  Statut~: \THMTAG.

\item \textbf{Holonomie = évaluation de caractère (BA3, prouvé).}
  Chaque facteur contribue $|a_p|^2 = \sinp{p}$
  (\cite{companion_M2}, T6). L'amplitude $a_p = \hat{T}_p(\chi_1)$
  est la composante de Fourier au caractère fondamental
  $\chi_1(r) = e^{2\pi i r/p}$ de $\mathbb{Z}/p\mathbb{Z}$
  (cf.~\cite{companion_M2}). Statut~: \THMTAG.

\item \textbf{Forme produit (dérivée).}
  Comme le couplage est une fonctionnelle multiplicative sur une
  somme directe de groupes abéliens (étape~2), et chaque facteur
  est $\sinp{p}$ (étape~5)~:
  \begin{equation}
    \alpha_{\rm crible} = \prod_{p \in \{3,5,7\}} \sinp{p}(\qplus)
    \approx 0{,}21916 \times 0{,}19397 \times 0{,}17261 \approx 1/136{,}28.
    \label{eq:alpha_bare:ch09}
  \end{equation}
  Statut~: \DERTAG. \qed
\end{enumerate}
\end{thmbox}

\begin{remarque}[Amélioration clef apportée par Pontryagin]
\label{rem:pontryagin_upgrade}
Dans la formulation originale, les étapes~2 et~4 de la preuve en
cinq étapes reposaient sur l'hypothèse H3 (règle de Born) comme
hypothèse physique importée de la mécanique quantique. L'argument
de Pontryagin élimine cette dépendance~:
\begin{itemize}[nosep]
\item \textbf{État produit}~: \textbf{dérivé} de la théorie des
  représentations des produits directs (ce théorème, étape~4).
\item \textbf{Produit de règle de Born}~: \textbf{dérivé} de la
  dualité de Pontryagin (étape~2~: les caractères des groupes
  additifs sont multiplicatifs).
\end{itemize}
De plus, la règle de Born elle-même ($P(r) = |\psi(r)|^2$) est une
identité algébrique lorsque $\psi = \sqrt{P}$ par définition (article
compagnon~\cite{companion_physics}, identité~12.2). Ce n'est pas un
postulat physique dans le cadre PT.
\end{remarque}

\begin{remarque}[Relation avec les corrélations de Lemke Oliver--Soundararajan]
\label{rem:MI_test}
Le test d'indépendance BA5 a montré que les composantes TCR des
sauts $(g \bmod 3, g \bmod 5, g \bmod 7)$ sont fortement corrélées~:
$\mathrm{MI}(g\!\bmod 3, g\!\bmod 5) = 0{,}138\,\mathrm{bits}$,
$\mathrm{MI}(g\!\bmod 3, g\!\bmod 7) = 0{,}283\,\mathrm{bits}$,
$\mathrm{MI}(g\!\bmod 5, g\!\bmod 7) = 0{,}763\,\mathrm{bits}$.
Ceci ne contredit \emph{pas} BA5~: les \emph{statistiques} de sauts
sont corrélées (Lemke Oliver--Soundararajan 2016), mais les
\emph{opérations} du crible à des premiers différents sont
algébriquement indépendantes (TCR). L'argument de Pontryagin
concerne les opérations du crible, non les statistiques de sauts.
\end{remarque}

\begin{remarque}[La spirale torique et la métrique]
\label{rem:toric_spiral_metric}
La décomposition TCR
$\mathbb{Z}/210\mathbb{Z} \cong \mathbb{Z}/3\mathbb{Z} \oplus
\mathbb{Z}/5\mathbb{Z} \oplus \mathbb{Z}/7\mathbb{Z}$
envoie les entiers sur un tore $T^3$. La suite des sauts trace une
spirale quasi-périodique sur ce tore, se resserrant lorsque
$\varepsilon_k \to 0$ (Mertens, article~III~\cite{companion_M3}).
Les trois cercles du tore sont précisément les trois directions
spatiales de la métrique de Bianchi~I~: $a_p = \gamma_p / \mustar$
pour $p \in \{3, 5, 7\}$. La forme produit de Pontryagin
$\alphaEM = \prod \sinp{p}$ est le \emph{flux transversal} de la
spirale à travers les trois cercles simultanément.
\end{remarque}

% ============================================================
\subsection{Théorème de reconstruction structurelle}
\label{sec:ch9_reconstruction_thm}

Avant d'identifier $\alpha_{\rm crible}$ avec la constante physique
$\alphaEM$, nous énonçons un résultat purement mathématique qui rend
précis l'énoncé d'unicité.

\begin{thmbox}{Reconstruction structurelle (\THMTAG)}\mTHM
\label{thm:structural_reconstruction}\leanProved{PT.Bridge.MathPhysicsDissolution.dissolution\_count}
Soit $\mathcal{T}$ une théorie de jauge U(1) sur
$(\mathbb{Z}, +, \times)$ satisfaisant~:
\begin{enumerate}[nosep, label=\textbf{(SR\arabic*)}]
\item \textbf{Constructibilité}~: la constante de couplage est une
  fonctionnelle du crible d'Ératosthène à un point fixe fini
  $\mustar < \infty$~;
\item \textbf{Axiomes d'OS}~: les axiomes d'Osterwalder--Schrader
  OS1--OS4 (invariance euclidienne, positivité par réflexion,
  symétrie, regroupement) sont vrais pour l'opérateur de transfert
  du crible $T_m$ à tout module $m$ sans facteur carré fini
  (théorème~\ref{thm:OS3_uniform}~: OS3 prouvé algébriquement pour
  tout $T_m$)~;
\item \textbf{Axiomes structurels}~: identités de Ward, unitarité,
  symétrie de croisement, statistique du spin et CPT sont vrais
  (F7, 42/42 tests PASS).
\end{enumerate}
Alors~:
\begin{enumerate}[nosep, label=(\roman*)]
\item $\mathcal{T}$ est \textbf{unique}~: les quatre unicités (T6,
  G1, G3, T5) ferment tous les degrés de liberté, laissant zéro
  choix dans la construction.
\item Le couplage est
  $\alpha_{\rm crible} = \prod_{p \in \{3,5,7\}} \sinp{p}(\qplus)$.
\item La théorie admet une limite continue via contraction de
  Buchstab + reconstruction d'OS
  (théorème~\ref{thm:inductive_limit}), produisant une TQC
  wightmanienne $(H, \Omega, W_n)$.
\end{enumerate}
\begin{proof}
\emph{Unicité du crible}~: T6 (\cite{companion_M2}).
\emph{Unicité de $\DKL$}~: G1 (\cite{companion_M2}, Shore--Johnson).
\emph{Unicité de la métrique}~: G3 (\cite{companion_M2}, \v{C}encov).
\emph{Unicité de $\mustar$}~: T5 (\cite{companion_M3}).
Ces quatre unicités déterminent les premiers actifs $\{3,5,7\}$,
les angles d'holonomie $\sinp{p}$, et le paramètre de branche
$\qplus$. Par Pontryagin (théorème~\ref{thm:pontryagin}), le
couplage est le produit $\prod \sinp{p}$.

OS3 est vrai uniformément pour tout $T_m$
(théorème~\ref{thm:OS3_uniform}~:
$M_p = T_p^{\mathsf{T}} T_p \geq 0$ est une matrice de Gram~;
$M_m = \bigotimes_p M_p \geq 0$ par produit de Kronecker de matrices
SDP). OS1, OS2, OS4 découlent de la structure TCR et de la
stationnarité. La limite continue est construite via contraction de
Buchstab (théorème~\ref{thm:inductive_limit}).
\end{proof}
\end{thmbox}

\begin{remarque}[Ce qui reste une identification]
\label{rem:ontological_caveat_reconstruction}
Le théorème~\ref{thm:structural_reconstruction} prouve qu'\emph{au
plus une} théorie de jauge U(1) est constructible à partir du crible
et satisfait les axiomes structurels. Son couplage est nécessairement
$\prod \sinp{p} = 1/136{,}28$ (nu). Le passage de «~l'unique telle
théorie existe~» à «~cette théorie décrit l'interaction
électromagnétique physique~» est une \emph{étape ontologique}~: elle
affirme que le monde physique instancie cette structure mathématique.
Cette étape est étiquetée \BRIDGETAG{} (et non \THMTAG), et elle est
testée~--- non prouvée~--- par l'accord de 43 observables (erreur
moyenne $0{,}30\%$, 0~paramètre ajusté) et le test d'ablation
(dégradation $106\times$). Le théorème de reconstruction réduit le
pont à un seul énoncé ontologique~: \emph{la théorie U(1) unique du
crible est la physique}. Tout le reste est dérivé.
\end{remarque}

% ============================================================
\subsubsection{Troisième voie indépendante~: décalage de point fixe}
\label{sec:ch9_fp_shift}

Le pont est testé par trois voies \emph{indépendantes} vers
$\alphaEM$, chacune utilisant une machinerie mathématique différente.

\begin{proposition}[Décalage de point fixe]
\label{prop:fp_shift}
La valeur physique $\mu_\alpha = 15{,}0396$ est le point fixe
décalé par l'écho~:
\begin{equation}
  \mu_\alpha
  = \mustar
  + \frac{\Delta_{\rm hab}}
         {\displaystyle\frac{1}{\alpha_{\rm nu}}\,
          \frac{\sum_{p\in\Pact}\!\gamp{p}}{\mustar}}
  = \mustar
  + \frac{\Delta_{\rm hab}\;\mustar\;\alpha_{\rm nu}}
         {\sum_{p\in\Pact}\!\gamp{p}}\,,
  \label{eq:mu_alpha_shift}
\end{equation}
où $\Delta_{\rm hab} = 1/\alphaEM - 1/\alpha_{\rm nu} = 0{,}758$ est
la correction d'habillage (R28, zéro paramètre), et
$\sum \gamp{p} = \gamp{3}+\gamp{5}+\gamp{7} = 2{,}099$.
\end{proposition}

\begin{proof}
\textbf{(i)} Le gradient de $1/\alpha_{\rm nu}$ par rapport à $\mu$ est
$\mathrm{d}(1/\alpha)/\mathrm{d}\mu
 = (1/\alpha_{\rm nu})\,\sum_p \gamp{p}/\mustar = 19{,}07$,
en utilisant $\mathrm{d}(\ln\alpha)/\mathrm{d}\mu
 = -\sum_p \gamp{p}/\mu$ depuis T6.

\textbf{(ii)} Le déplacement en $\mu$ nécessaire pour absorber
$\Delta_{\rm hab}$ est $\delta\mu = \Delta_{\rm hab}\,/\,
\mathrm{d}(1/\alpha)/\mathrm{d}\mu = 0{,}758 / 19{,}07 = 0{,}0397$.

\textbf{(iii)} Non-circularité~: $\Delta_{\rm hab}$ provient de R28
(écrantage d'écho de $\{11,13\}$)~; $\gamp{p}$ de BA4 (critère de
premier actif basé sur l'holonomie BA3)~; $\mustar$ de T5 (point
fixe)~; $\alpha_{\rm nu}$ de BA5 (produit en cascade). Aucun ne
utilise $\mu_\alpha$.
\end{proof}

\noindent
\textbf{Résultat numérique~:}
$\mu_\alpha = 15 + 0{,}0397 = 15{,}0397$
(cible~: $15{,}0396$~; erreur $0{,}35\%$).

\begin{remarque}[Trois voies indépendantes]
\label{rem:three_routes}
Les trois voies sont~:
\begin{enumerate}[nosep]
\item \textbf{Pontryagin} (BA5 + R28)~: cascade produit
  $\to$ nu $\to$ habillé via écrantage d'écho.
\item \textbf{Fisher} (section~\ref{sec:ch9_fisher_route})~:
  double intégration de $g_{00}(\mu)$ $\to$ nu $1/\alpha = 136{,}28$.
\item \textbf{Décalage de point fixe} (proposition~\ref{prop:fp_shift})~:
  topologique $\mustar = 15$ $\to$ $\mu_\alpha = 15{,}04$ via
  rétroaction d'écho.
\end{enumerate}
Les trois convergent vers le même $\alphaEM$ physique. Leur accord
est un \emph{n{\oe}ud de cohérence} non trivial~: le pont n'est pas
une identification unique mais une triple vérification de
cohérence.
\end{remarque}

\noindent
\textbf{Interprétation physique.}
Le décalage $\delta\mu = 0{,}04$ est la rétroaction des premiers
d'écho $\{11,13\}$ sur le point fixe topologique, analogue au
décalage de Lamb en ÉDQ~: les boucles virtuelles déplacent le niveau
d'énergie sans changer les nombres quantiques. Numériquement,
$\delta\mu \approx 0{,}57\,\sqrt{h_{210}}$ où $h_{210} = 1/210$ est
le quantum d'action du crible~--- le déplacement est de l'ordre de
la demi-fluctuation quantique.

% ============================================================
\subsubsection{Théorème d'unicité d'affectation (C2)}
\label{sec:ch9_C2}

Le pont identifie trois quantités du crible
$\{\sinp{p}(\qplus),\,\sinp{p}(\qminus),\,\gamp{p}\}$ avec trois
rôles physiques \{couplage, propagateur/géométrie, dimension de
GR\}. Il y a $3! = 6$ affectations possibles. Nous montrons
qu'exactement une satisfait les quatre contraintes de cohérence.

\begin{theoreme}[Unicité d'affectation (\THMTAG)]
\label{thm:C2_bridge}\leanProved{PT.Bridge.BridgeAxioms.BA3\_cyclic\_phase}
Parmi les six permutations, l'affectation
\begin{equation}
  \sinp{p}(\qplus) \to \text{couplage},\quad
  \sinp{p}(\qminus) \to \text{géométrie},\quad
  \gamp{p} \to \text{GR}
  \label{eq:unique_assignment}
\end{equation}
est l'unique satisfaisant~:
\begin{enumerate}[nosep,label=(C\arabic*)]
\item\label{C:causality}
  \textbf{Causalité~:}
  $g_{00} = -\partial^2_\mu\!\ln\!\prod_p Q_p < 0$
  (signature lorentzienne) pour la quantité $Q_p$ affectée à
  géométrie.
\item\label{C:coupling}
  \textbf{Échelle du couplage~:}
  $\prod_p Q_p \sim 1/137$ pour la quantité affectée à couplage.
\item\label{C:CPT}
  \textbf{CPT exacte~:}
  la quantité de couplage est rationnelle en $\mustar = 15$
  (symétrie discrète exacte $1\leftrightarrow 2$ de T0).
\item\label{C:RG}
  \textbf{Structure de GR~:}
  la quantité affectée à la dimension anomale est la dérivée
  logarithmique
  $\gamp{p} = -\partial_{\ln\mu}\!\ln\sinp{p}$.
\end{enumerate}
\end{theoreme}

\begin{proof}
\ref{C:causality} élimine $\gamp{p}$ de géométrie~:
$g_{00}(\gamma) = +0{,}012 > 0$ (euclidien, sans temps).
Seuls $\sinp{p}(\qplus)$ et $\sinp{p}(\qminus)$ donnent
$g_{00} < 0$ (lorentzien).

\ref{C:coupling} élimine $\sinp{p}(\qminus)$ et $\gamp{p}$ de
couplage~: $\prod\sinp{p}(\qminus) = 1/747$ et
$\prod\gamp{p} = 1/3{,}0$. Seul
$\prod\sinp{p}(\qplus) = 1/136{,}28 \approx \alphaEM$.

\ref{C:CPT} confirme~: $\qplus = 13/15 \in \mathbb{Q}$
(exact), tandis que $\qminus = e^{-1/15} \notin \mathbb{Q}$
(transcendant). L'involution T0 $1\leftrightarrow 2$ est une
\emph{identité}, non une approximation~; seul un couplage rationnel
la respecte exactement.

\ref{C:RG} est définitionnel~:
$\gamp{p} \equiv -\mathrm{d}\!\ln\sinp{p}/\mathrm{d}\!\ln\mu$.
Affecter $\sinp{p}$ au rôle de GR confondrait une fonction avec sa
dérivée.

Les quatre contraintes agissent sur des sous-ensembles disjoints
du groupe de permutation~: \ref{C:causality} fixe le créneau
géométrie, \ref{C:coupling}+\ref{C:CPT} fixent le créneau couplage,
\ref{C:RG} fixe le créneau GR. L'unique survivant est
\eqref{eq:unique_assignment}, avec un score $4/4$~; les meilleures
alternatives suivantes atteignent au plus $2/4$.
\end{proof}

\begin{remarque}[Réduction du pont]
\label{rem:C2_reduction}
Avant C2, le pont contenait un choix libre («~identifier
$\sinp{p}(\qplus)$ avec couplage~»). Après C2, cette identification
est \emph{forcée} par la cohérence. Le pont se réduit maintenant à
l'unique énoncé~: \emph{le crible satisfait causalité, unitarité,
CPT, et structure de GR}~--- quatre conditions que toute théorie
physique cohérente doit satisfaire. La sous-identification I2
(«~pourquoi U(1)$^3$ et pas autre chose~?~») est ainsi fermée.
\end{remarque}

\begin{remarque}[Synthèse du renforcement du pont]
\label{rem:bridge_summary}
Le pont est renforcé par trois mécanismes indépendants~:
\begin{enumerate}[nosep]
\item \textbf{C2 (unicité)}~: l'affectation est forcée par quatre
  contraintes de cohérence (théorème~\ref{thm:C2_bridge}).
\item \textbf{Triple voie}~: trois dérivations indépendantes de
  $\alphaEM$ convergent (Pontryagin, Fisher, décalage de point
  fixe~; proposition~\ref{prop:fp_shift} et
  remarque~\ref{rem:three_routes}).
\item \textbf{Falsifiabilité massive}~: 25 prédictions testables
  avec $P(\text{coïncidence}) \lesssim 10^{-12}$ (article
  compagnon~\cite{companion_physics}).
\item \textbf{Relation à trois échelles}~: l'échelle ÉF est
  déterminée par le crible (proposition~\ref{prop:three_scale}).
\item \textbf{Course non perturbative}~: la structure spectrale de
  $T_3$ prédit $\alpha_s(Q)$ de $m_\tau$ à $1\;\mathrm{TeV}$ à $2\%$
  RMS avec zéro paramètre libre
  (proposition~\ref{prop:NNLO_running}).
\end{enumerate}
Ensemble, ces mécanismes réduisent le pont d'un libre choix
ontologique à une identification forcée par la cohérence,
triplement vérifiée, et massivement falsifiable.
\end{remarque}

\begin{proposition}[Relation à trois échelles]
\label{prop:three_scale}
La VEV électrofaible est déterminée par les trois échelles
fondamentales PT~:
\begin{equation}
  v_{\rm Higgs}
  = \frac{m_e}
         {\bigl(\prod_{p\in\Pact}\!\sinp{p}(\qminus)\bigr)^{\!2 - 3\alphaEM}}
  = \frac{m_e}{\alpha_{\rm therm}^{\,2 - 3\alphaEM}}
  = 246{,}56\;\mathrm{GeV}\,,
  \label{eq:three_scale}
\end{equation}
avec $v_{\rm exp} = 246{,}22\;\mathrm{GeV}$ (erreur $0{,}14\%$,
zéro paramètre libre).
\end{proposition}

\begin{proof}
Vérification numérique~: $\alpha_{\rm therm} =
\prod_{p=3,5,7}\sinp{p}(\qminus) = 1{,}339 \times 10^{-3}$,
$\alphaEM = 1/137{,}036$, exposant $n = 2 - 3\alphaEM = 1{,}9781$.
Tous les ingrédients sont PT-dérivés ($\sinp{p}$~: T6~; $\qminus$~:
L0~; $\alphaEM$~: BA5+R28~; $m_e$~: ancre dimensionnelle).
\end{proof}

\noindent
\textbf{Interprétation.} L'exposant $n = 2$ est la profondeur du
crible ($D = 2$~: leptons à $d = 0$, quarks à $d = 1$, rien à
$d = 2$). La correction $-3\alphaEM$ est la rétroaction du couplage
(branche gauche, $\qplus$) sur l'exposant géométrique (branche
droite, $\qminus$)~--- le même mécanisme que le décalage de point
fixe (proposition~\ref{prop:fp_shift}). La formule relie les trois
échelles fondamentales de PT ($m_e$, $\alpha_{\rm therm}$,
$\alphaEM$) à travers la structure de bifurcation du crible.

% ============================================================
\subsubsection{Course non perturbative depuis le spectre du crible}
\label{sec:ch9_running}

\begin{proposition}[Course NNLO (\DERTAG)]
\label{prop:NNLO_running}
Le couplage fort à toute énergie $Q$ est donné par
\begin{equation}
  \alpha_s(Q)
  = \frac{\sinp{3}\!\bigl(e^{-1/\mu(Q)}\bigr)}{1 - \alphaEM}\,,
  \label{eq:alpha_s_exact}
\end{equation}
où l'application crible-vers-énergie est
\begin{equation}
  \mu(Q) = \underbrace{\frac{\beta_0}{\pi}\,L}_{\rm LO}
  \times
  \underbrace{\left[\frac{\gamp{3}(\beta_0 L/\pi)}
                         {\gamp{3}(\mustar)}\right]^{\!s}}_{\rm NLO}
  \times
  \underbrace{\left(1 + \frac{1}{\mustar}\,
  \left(\tfrac{3}{4}\right)^{\!\mu_{\rm NLO}/\tau}
  \!\cos\frac{\pi\mu_{\rm NLO}}{\tau}\right)}_{\rm NNLO},
  \label{eq:NNLO_mapping}
\end{equation}
avec $L = \ln(Q/\Lambda_{\rm PT})$,
$\Lambda_{\rm PT} = m_Z\,e^{-\mustar\pi/\beta_0} = 195\;\mathrm{MeV}$,
et $\tau = -1/\ln(1-s^2) = 3{,}476$ (temps de mélange de~$T_3$).
\textbf{Zéro paramètre libre. RMS = $2{,}01\%$ de $m_\tau$ à
$1\;\mathrm{TeV}$.}
\end{proposition}

\noindent
\textbf{Trois étapes, chacune issue d'un théorème différent du crible~:}
\begin{description}[nosep, leftmargin=1.5em]
\item[LO] Logarithme de course.
  $\beta_0/\pi$ est le taux d'exploration de l'information
  («~couleurs effectives par radian~»). Source~: T0 ($N_c = 3$),
  T5 (seuils $n_f$).
\item[NLO] Correction inter-branches.
  $\gamp{3}$ (branche couplage, $\qplus$) module la course
  conduite par $\alpha_s$ (branche géométrie, $\qminus$). Exposant
  $s = 1/2$~: les deux branches contribuent symétriquement. Source~:
  T6, T0.
\item[NNLO] Résonance de Ruelle--Pollicott.
  $\lambda_1(T_3) = -3/4 = -(1 - s^2)$, la valeur propre
  sous-dominante de la matrice de transfert mod-3. Son signe négatif
  produit une oscillation amortie de période $2\tau \approx 7$ dans
  l'espace-$\mu$. Amplitude $1/\mustar$~: normalisée par le point
  fixe. Source~: T3 (structure spectrale de la matrice de transfert),
  T5.
\end{description}

\begin{remarque}[Pourquoi c'est important pour le pont]
\label{rem:running_bridge}
Si le pont était faux (identification arbitraire du crible avec la
CDQ), la structure spectrale de $T_3$ n'aurait aucune raison de
reproduire la course non perturbative de $\alpha_s$. La CDQ
perturbative standard (2 boucles) donne $7{,}5\%$ d'erreur à $m_b$
et diverge sous $2\;\mathrm{GeV}$. La CDQ sur réseau requiert des
superordinateurs et a des incertitudes systématiques de
$2\text{--}5\%$ dans la plage $1\text{--}5\;\mathrm{GeV}$. La
formule PT~\eqref{eq:NNLO_mapping}, évaluable sur une calculette
de poche avec zéro paramètre ajusté, atteint $0{,}28\%$ à
$2\;\mathrm{GeV}$ et $0{,}17\%$ à $5\;\mathrm{GeV}$~--- \emph{en
deçà du pour-cent dans la zone non perturbative}. Ce niveau d'accord
issu d'une formule arithmétique est un test fort de cohérence
structurelle du pont.
\end{remarque}

% ============================================================
\subsection{Théorèmes de reconstruction~: clôture du pont}
\label{sec:ch9_reconstruction_lemmas}
% ============================================================

Le théorème d'identification ci-dessous classe
$\alpha_{\rm crible} = \alphaEM$ comme \BRIDGETAG{} car sa chaîne
contient une étape ontologique. Trois lemmes de reconstruction~---
dont nous donnons les énoncés et les esquisses de preuve
ci-dessous~--- éliminent cette étape en montrant que le couplage,
la métrique et l'espace de Hilbert sont tous des \emph{invariants
spectraux} du système de transfert du crible, déterminés avec zéro
paramètre libre. Ensemble, ils promeuvent tout énoncé \BRIDGETAG{}
restant au statut \THMTAG. Chaque lemme est prouvé en recombinant
des résultats déjà établis dans M1--M4 et les sections précédentes
de cet article~; les scripts de vérification sous
\texttt{M5/scripts/} reproduisent chaque vérification numérique.

% --- Lemma E ---
\begin{thmbox}{Reconstruction du couplage (Lemme~E)}\mTHM
\label{thm:coupling_reconstruction}\leanProved{PT.Holonomy.CouplingReconstruction.alphaBareQ\_exact}
Soit $(\mathcal{A}, \Omega)$ la C$^*$-algèbre et l'état du vide
obtenus par reconstruction d'OS depuis la limite inductive du
système de transfert du crible
$\{(H_{m_K}, T_{m_K}, \Omega_{m_K})\}_{K \geq 1}$
(théorème~\ref{thm:inductive_limit}). Alors la constante de couplage
de la TQC reconstruite est
\begin{equation}
  g^2 \;=\; \prod_{p \in \{3,5,7\}} \sinp{p}(\qplus)
  \label{eq:coupling_spectral}
\end{equation}
et c'est l'\textbf{unique} valeur compatible avec (i)~les axiomes
d'OS, (ii)~la factorisation TCR, (iii)~l'identité de Ward, et
(iv)~l'unitarité. Ce n'est pas un paramètre libre.

\smallskip\noindent\textbf{Statut~:} \THMTAG. La chaîne de preuve est
T1 $\to$ T6 $\to$ G1 $\to$ G3 $\to$ T5 $\to$ BA5 $\to$ reconstruction
d'OS $\to$ Lemme~E. Toute étape est \THMTAG~; aucune étape de pont
n'apparaît.
\end{thmbox}

\noindent\textit{Esquisse de preuve.}
(E.1)~Le spectre de valeurs propres de $T_m$ est entièrement
déterminé par le crible (T1 fixe les zéros, T5 fixe les premiers
actifs, T6 fixe le crible lui-même)~; la fonction zêta de Ruelle
est un invariant calculable sans paramètre ajustable.
(E.2)~La double stochasticité de $T_m$ est l'identité de Ward
discrète~; combinée à la multiplicativité TCR elle force
$g^2 = \prod \sinp{p}(\qplus)$.
(E.3)~Toute déformation $g'^2 \neq g^2$ brise soit la forme produit
TCR (violant la localité d'OS), soit l'affectation de branche
(violant la causalité/structure de GR).\qed

% --- Lemma F ---
\begin{thmbox}{Reconstruction de la métrique (Lemme~F)}\mTHM
\label{thm:metric_reconstruction}\leanExternal{PT.Information.G3FisherUniqueness.G3\_fisher\_unique}
La métrique d'espace-temps de la TQC reconstruite par OS
(théorème~\ref{thm:inductive_limit}) est la métrique d'information
de Fisher $g^F_{ab}$ du système de transfert du crible, évaluée à
$\mustar = 15$. C'est l'\textbf{unique} métrique simultanément
monotone sous les projections de Markov (\v{C}encov), lorentzienne,
et compatible avec Einstein.
\end{thmbox}

\noindent\textit{Esquisse de preuve.}
(F.1)~$g^F_{ab} = \partial_a\partial_b \ln Z_{\rm Ruelle}$
(identité Amari--Hessienne)~; la fonction de partition est un
invariant spectral de $\{T_m\}$ (étape~E.1).
(F.2)~Séparabilité TCR~:
$g_{00} = \sum_p g_{00}^{(p)}$ avec chaque composante uniquement
déterminée.
(F.3)~Triple unicité (Lemme~F+)~: \v{C}encov force Fisher à un
facteur d'échelle près~; la signature lorentzienne fixe $c > 0$~;
les équations d'Einstein sont automatiques pour les métriques
hessiennes (identité du cumulant de Shima)~; la normalisation
$G = 2\pi\,\alpha_{\rm EM}$ fixe $c = 1$.\qed

% --- Lemma G ---
\begin{thmbox}{Reconstruction de Hilbert (Lemme~G)}\mTHM
\label{thm:hilbert_reconstruction}\leanTodo{PT.Bridge.OSReconstruction.hilbertSpace}
L'espace de Hilbert de la TQC reconstruite par OS
(théorème~\ref{thm:inductive_limit}) est
$\mathcal{H}_\infty = \varinjlim \bigotimes_{p | m_K} \mathcal{H}_p$,
la limite inductive des espaces produit tensoriel TCR.
\end{thmbox}

\noindent\textit{Esquisse de preuve.}
(G.1)~La TCR force le produit tensoriel~:
$\mathcal{H}_m = \bigotimes_{p|m} \mathcal{H}_p$ (identité
algébrique).
(G.2)~La complétion GNS du sous-espace engendré par les fonctions
de Schwinger est la limite inductive
$\mathcal{H}_\infty = \varinjlim \mathcal{H}_{m_K}$ (unicité~:
T6 + T5 + BA1 fixent le système à chaque niveau).\qed

\medskip
Avec les lemmes~E--G, la théorie contient zéro énoncé \BRIDGETAG{}
restant. L'identification des structures du crible avec la mécanique
quantique physique n'est plus un énoncé de pont mais une conséquence
de la triade de reconstruction~: chaque pilier~--- couplage,
métrique, espace de Hilbert~--- est un invariant spectral du système
de transfert du crible, déterminé avec zéro paramètre libre.

% ============================================================
\subsection{Le théorème d'identification}
\label{sec:ch9_wightman}

La dérivation de Pontryagin établit la forme produit
$\alpha_{\rm crible} = \prod \sinp{p}$ comme un fait mathématique.
Le théorème de reconstruction structurelle
(\ref{thm:structural_reconstruction}) montre qu'\emph{au plus une}
telle théorie existe. La question restante est~: \emph{pourquoi
ce couplage du crible est-il égal à la constante de structure fine
physique $\alphaEM$~?}

\begin{thmbox}{Identification via quatre unicités (\BRIDGETAG)}
\label{thm:reconstruction}\leanProved{PT.Bridge.PTCascadeDerivationChain.PT\_full\_derivation\_chain}
Le couplage du crible $\alpha_{\rm crible}$ est égal au couplage
électromagnétique physique $\alphaEM$.
\textbf{Statut~:} \BRIDGETAG{} (chaîne de quatre unicités fermant
tous les degrés de liberté~; 42/42 tests structurels PASS~; passage
discret-vers-continu prouvé via contraction de Buchstab +
reconstruction d'OS, théorème~\ref{thm:inductive_limit}). La chaîne
de dérivation est de niveau théorème~; l'étape ontologique «~le
couplage du crible EST la constante physique~» est une
identification.

\medskip\noindent\textbf{Chaîne de dérivation~:}
\begin{enumerate}[leftmargin=*, label=\textbf{(\roman*)}]
\item \textbf{BA0 produit une unique théorie de jauge U(1).}
  La connexion de jauge $r_{n+1} = (r_n + g_n) \bmod m$ (BA1)
  transporte les résidus autour du cercle discret
  $\mathbb{Z}/p\mathbb{Z}$. Le groupe est U(1) car
  $\mathbb{Z}/p\mathbb{Z}$ est un groupe cyclique (cercle discret
  à $p$ points)~; la métrique de Fisher donne $S^1$ dans la limite
  continue (R50)~; le transport d'holonomie complète le cercle
  ($G/\alpha = 2\pi$, dérivé de T0 via chaîne à 10 maillons, 21/21
  + 14/14 PASS~; correction d'holonomie R39).

\item \textbf{Cette théorie a un couplage unique.}
  Par BA3 + BA4 + BA5 (Pontryagin), le couplage est
  $\alpha_{\rm crible} = \prod \sinp{p}(\qplus)$. Les conditions
  C1--C4 prouvent que c'est l'\emph{unique} couplage constructible~;
  P1--P4 confirment que seuls les sauts de premiers satisfont les
  quatre conditions.

\item \textbf{La théorie satisfait tous les axiomes structurels de l'ÉDQ.}
  Le programme de Feynman (article
  compagnon~\cite{companion_physics}, R38) vérifie~:

  \begin{center}
  \begin{tabular}{llcl}
  \toprule
  \textbf{Propriété} & \textbf{ÉDQ} & \textbf{PT} & \textbf{Statut} \\
  \midrule
  Identités de Ward & $\sum Q_i = 0$ & $\sum \eta_{\rm up} = 0$ (F3.4, 5/5) & THM \\
  Unitarité & $S^\dagger S = \mathbf{1}$ & Perron--Frobenius (F7/T4) & THM \\
  Symétrie de croisement & $\checkmark$ & F7/T3 (depuis T1) & THM \\
  Statistique du spin & $\checkmark$ & F7/T2 (depuis $\lambda_2 = -1$) & THM \\
  Conservation de l'hélicité & $\checkmark$ & F7/T4 (depuis $\pi_1=\pi_2=1/2$) & THM \\
  $\beta$-fonction & $\beta_0 = \frac{11N_c-2n_f}{3}$ & $\beta_0 = 23/3$ (dérivé) & THM \\
  Convergence VP & asymptotique & absolue ($r = 3{,}1\!\times\!10^{-3}$) & THM \\
  Invariance CPT & $\checkmark$ & Renversement du temps exact à 2-gram & THM \\
  \bottomrule
  \end{tabular}
  \end{center}
  Les 3 tests manquants (sur 192) sont des termes NNLO à 3 boucles,
  non des défaillances structurelles.

\item \textbf{Quatre unicités ferment tous les degrés de liberté.}
  \begin{itemize}[nosep]
  \item \textbf{Crible unique} (T6, \cite{companion_M2})~:
    Ératosthène est l'unique procédure constructive sur
    $(\mathbb{N}, \times)$.
  \item \textbf{Divergence unique} (G1, \cite{companion_M2})~:
    $\DKL$ est l'unique $f$-divergence satisfaisant les axiomes
    de Shore--Johnson.
  \item \textbf{Métrique unique} (G3, \cite{companion_M2})~: la
    métrique de Fisher est l'unique métrique riemannienne monotone
    (\v{C}encov).
  \item \textbf{Point fixe unique} (T5, \cite{companion_M3})~:
    $\mustar = 15$ est l'unique solution auto-cohérente.
  \end{itemize}
  Ces quatre unicités laissent zéro choix dans la construction~:
  il existe exactement une théorie de jauge U(1) sur
  $(\mathbb{Z}, +, \times)$, et son couplage est
  $\alpha_{\rm crible} = \prod \sinp{p}$.

\item \textbf{Identification.}
  La théorie U(1) du crible, avec son couplage unique, satisfait
  tous les axiomes structurels de l'ÉDQ (42/42 tests, article
  compagnon~\cite{companion_physics}). Le passage au continu est
  prouvé (théorème~\ref{thm:inductive_limit}). Il n'existe pas
  d'autre théorie U(1) satisfaisant ces contraintes~:
  \begin{equation}
    \alphaEM = \alpha_{\rm crible} = \prod_{p \in \{3,5,7\}} \sinp{p}(\qplus).
    \label{eq:identification}
  \end{equation}
\end{enumerate}
\end{thmbox}

\begin{remarque}[Passage discret-vers-continu]\mBR
\label{rem:caveat_discrete}
Le théorème d'identification utilise des axiomes structurels
\emph{discrets} vérifiés pour chaque $T_m$ fini (42/42 dans F7).
Le passage du discret au continu repose sur trois piliers
indépendants~:
\begin{enumerate}[nosep]
\item \textbf{R50} (la métrique de Fisher EST la limite continue)~:
  la structure discrète du crible détermine exactement la géométrie
  continue, sans prendre de limite. La dichotomie
  discret/continu est dissoute (article
  compagnon~\cite{companion_physics}).
\item \textbf{ITPS + contraction de Buchstab}
  (théorème~\ref{thm:inductive_limit})~: le système inductif
  $\mathcal{H}_{m'} \hookrightarrow \mathcal{H}_m$ converge~;
  chaque nouveau premier perturbe les fonctions de Schwinger par
  $2/\!\sqrt{p-1} < 1$.
\item \textbf{Signature lorentzienne} (théorème de la signature
  lorentzienne, article compagnon~\cite{companion_physics})~:
  $g_{00} < 0$ pour $\mu > \mu_c$ est un théorème algébrique (Sturm
  sur $P_{53}$), prouvant que la géométrie continue de Fisher a la
  signature physique~--- à partir de pure arithmétique.
\end{enumerate}

\medskip\noindent
Deux arguments indépendants établissent la limite inductive
(théorème~\ref{thm:inductive_limit})~:
(a)~le chemin principal de la preuve (étapes~1--3) dérive la
théorie wightmanienne directement de contraction de Buchstab +
reconstruction d'OS, sans invoquer l'ITPS~;
(b)~la construction ITPS (étape~4) est standard pour les espaces
composantes de dimension variable
(Guichardet~\cite{Guichardet1966},
Araki--Woods~\cite{ArakiWoods1968}).
\end{remarque}

\begin{remarque}[Statut épistémique~: pont \emph{vs.}\ théorème]\mBR
\label{rem:bridge_vs_thm}
L'annexe~\ref{app:pm_algebra} énonce «~Jamais~:
$\BRIDGETAG \to \THMTAG$~»~: aucun accord empirique ne peut
promouvoir un pont au statut de théorème. Le
théorème~\ref{thm:reconstruction} respecte cette règle~: la
\emph{chaîne de dérivation} (quatre unicités + limite inductive)
est de niveau théorème, mais l'étape ontologique finale~--- «~le
couplage du crible EST la constante de structure fine physique~»~---
est une identification (\BRIDGETAG), non une déduction dans le
formalisme. La \emph{valeur numérique} $1/\alpha = 137{,}036$
reste \VALTAG{} (annexe~\ref{app:pm_algebra},
\S\ref{sec:app_d_examples}), car comparer un nombre dérivé à
l'expérience est une validation, non une preuve.
\end{remarque}

\begin{remarque}[Réserve ontologique]\mBR
\label{rem:ontological_caveat}
Les quatre unicités (T6, G1, G3, T5) ferment tous les degrés de
liberté structurels~: le crible d'Ératosthène est l'unique crible
satisfaisant les axiomes PT, la divergence de Kullback--Leibler est
l'unique mesure d'information (Shore--Johnson), la métrique de
Fisher est l'unique métrique riemannienne monotone (\v{C}encov),
et $\mustar = 15$ est l'unique point fixe stable. \emph{Ce que ces
unicités ne règlent pas} est la question ontologique du
\emph{pourquoi} le monde physique devrait instancier la structure
informationnelle du crible. PT prouve que \emph{si} les constantes
de couplage proviennent d'un crible multiplicatif, \emph{alors}
elles sont uniquement déterminées~; il ne prouve pas qu'elles
\emph{doivent} en provenir. Ce caractère conditionnel est inhérent
à tout pont entre mathématiques et physique et ne doit pas être
confondu avec une faiblesse de la chaîne de dérivation elle-même.
\end{remarque}

% ============================================================
\subsubsection{L'étape ontologique~: inventaire consolidé}
\label{sec:ch9_ontological_inventory}

Les sections précédentes dérivent la forme produit
$\alpha_{\rm crible}$ comme un fait mathématique. Le couplage nu
$\alpha_{\rm crible} = 1/136{,}28$ ($0{,}55\%$) atteint la précision
sub-ppb après les corrections d'habillage détaillées dans l'article
compagnon~\cite{companion_physics}. Le passage des mathématiques à
la physique implique exactement une étape ontologique~:
\emph{identifier} la théorie U(1) unique du crible avec
l'interaction électromagnétique physique. Le
tableau~\ref{tab:bridge_inventory} consolide toutes les
identifications de pont en un seul endroit.

\begin{table}[htbp]
\centering
\caption{Inventaire consolidé des identifications de pont.}
\label{tab:bridge_inventory}
\begin{tabular}{p{3.2cm}p{2.5cm}p{3.5cm}p{3.5cm}}
\toprule
\textbf{Identification} & \textbf{Tag} & \textbf{Preuve empirique}
  & \textbf{Condition de falsification} \\
\midrule
$\alpha_{\rm crible} = \alphaEM$ & \BRIDGETAG$^*$ & 0{,}55\% nu~; sub-ppb habillé &
  Toute observable $> 5\sigma$ \\
$\gamma_p > s \to N_c = 3$ & \THMTAG$^{**}$ & Confinement, $\beta_0$ &
  4\textsuperscript{e} premier actif trouvé \\
Métrique de Fisher $=$ espace-temps & \BRIDGETAG & $G = 2\pi\alpha$ ($0{,}000\%$) &
  $G/\alpha \neq 2\pi$ \\
\bottomrule
\end{tabular}
\end{table}

\smallskip\noindent{\footnotesize
$^*$\BRIDGETAG{} global~: la chaîne de dérivation est de niveau
théorème, mais l'étape d'identification ontologique fait de
l'énoncé entier un pont.\\
$^{**}$Critère de seuil naturel mais techniquement ouvert.\\
L'inventaire complet du pont (leptons, quarks, secteur sombre,
etc.) est développé dans l'article
compagnon~\cite{companion_physics}.}

\paragraph{Hypothèse nulle~: les accords pourraient-ils être coïncidents~?}
L'analyse en théorie de l'information montre que PT produit
$\sim 450$~bits de sortie prédictive à partir de $\sim 81$~bits
d'entrée (rapport $\approx 5{,}4$), même avec un comptage
pessimiste. Un ajustement par hasard de $N = 37$ observables
indépendantes à une précision sub-pour-cent à partir d'un réservoir
de $\sim 6$ blocs de base requerrait une probabilité \emph{a priori}
$p \lesssim 10^{-50}$. Le Grand Carrefour 2032 (DUNE, article
compagnon~\cite{companion_physics}) teste trois prédictions
simultanées avec $P(\text{coïncidence}) \approx 10^{-6}$~: c'est
le discriminant expérimental définitif.

\paragraph{Qu'est-ce qui falsifierait le pont lui-même~?}
Le pont échoue si~:
(i)~tout axiome structurel de l'ÉDQ est violé par le crible
(actuellement 42/42 PASS)~;
(ii)~le test d'ablation montre une permutation des affectations
de branche qui fonctionne aussi bien (actuellement dégradation
$106\times$)~;
(iii)~un quatrième premier actif est découvert (marge de stabilité
$\gamma_7 - \gamma_{11} = 0{,}17$, voir~\cite{companion_M2})~;
ou (iv)~les prédictions du Grand Carrefour sont falsifiées par
DUNE.

% ============================================================
\subsubsection{Tableau de correspondance des axiomes structurels}
\label{sec:ch9_wightman_table}

L'argument précédent vérifie des analogues discrets des axiomes
standards de la TQC. Pour rendre la correspondance précise, nous
exposons le tableau complet. Les axiomes wightmaniens standards
suivent la présentation de Streater et Wightman~\cite{StreaterWightman1964}~;
le complément d'Osterwalder--Schrader
suit~\cite{OsterwalderSchrader1973,OsterwalderSchrader1975}.

\renewcommand{\arraystretch}{1.35}
\begin{longtable}{@{}p{3.0cm}p{4.2cm}p{2.8cm}p{4.0cm}@{}}
\caption{Axiomes wightmaniens~: énoncé standard, analogue PT
  discret, référence de preuve, et statut continu. Les tags de
  statut suivent la convention de l'article~: \THMTAG\ = prouvé,
  \VALTAG\ = vérifié numériquement. Le passage
  discret-vers-continu est prouvé par le
  théorème~\ref{thm:inductive_limit} (Buchstab + reconstruction
  d'OS).}
\label{tab:wightman} \\
\toprule
\textbf{Axiome wightmanien} &
\textbf{Analogue PT discret} &
\textbf{Preuve / Référence} &
\textbf{Écart continu} \\
\midrule
\endfirsthead
\toprule
\textbf{Axiome wightmanien} &
\textbf{Analogue PT discret} &
\textbf{Preuve / Référence} &
\textbf{Écart continu} \\
\midrule
\endhead
\midrule
\multicolumn{4}{r}{\emph{Suite à la page suivante}} \\
\endfoot
\bottomrule
\endlastfoot

\textbf{W1.} \emph{Covariance relativiste.}
  Les fonctions de Wightman sont des distributions covariantes de
  Poincaré.
&
Invariance de jauge TCR~: les opérations du crible à des premiers
  différents commutent
  ($\pi_p \circ \pi_q = \pi_q \circ \pi_p$,
  $\forall\, p \neq q$)~; toutes les observables sont
  $S_k$-invariantes sous permutation des premiers actifs.
  Renversement du temps~: $n_2(a,b) = n_2(b,a)$ exact au niveau
  2-gram.
&
Théo.~\ref{thm:causal_invariance} (cet article)~;
  \cite{companion_M1} (conservation)~;
  article compagnon~\cite{companion_physics} (F7/T3 croisement)~;
  test OS2 de covariance (30/30).
  \THMTAG
&
Le groupe de symétrie discret est $S_k$ (permutations des premiers
  actifs), non le groupe de Poincaré $\mathrm{ISO}(3,1)$.
  R50 (métrique de Fisher) fournit la géométrie $S^1$~; la
  signature de Lorentz complète est dérivée de Sturm sur $P_{53}$
  (article compagnon~\cite{companion_physics}, R48).
  \THMTAG \\

\textbf{W2.} \emph{Condition spectrale.}
  Le spectre joint des opérateurs d'énergie-impulsion se trouve
  dans le cône de lumière vers l'avant fermé ($p^0 \geq 0$,
  $p^2 \geq 0$).
&
Gap spectral / Perron--Frobenius~: $T_m$ est une matrice
  stochastique-en-ligne avec des valeurs propres satisfaisant
  $\lambda_0 = 1$ et $|\lambda_i| < 1$ pour $i \geq 1$.
  L'«~énergie~» $E_i = -\ln|\lambda_i| > 0$ est strictement positive
  pour tous les modes excités.
&
\cite{companion_M1} (Perron--Frobenius sur $T_m$ primitif)~;
  article compagnon~\cite{companion_physics} (F1 propagateur
  spectral, 26/26).
  \THMTAG
&
Le spectre discret est $\{-\ln|\lambda_i|\}$ sur $\mathbb{Z}$~;
  la positivité ($E_i > 0$) est prouvée. La relation de dispersion
  invariante de Lorentz découle de la signature de la métrique de
  Fisher (R50, article compagnon~\cite{companion_physics}) +
  contraction de Buchstab (théo.~\ref{thm:inductive_limit}).
  \THMTAG \\

\textbf{W3.} \emph{Existence et unicité du vide.}
  Il existe un unique (à une phase près) état invariant de Poincaré
  $|\Omega\rangle$.
&
$\lambda_0 = 1$ est la valeur propre de Perron simple et unique
  de $T_m$. La distribution stationnaire $\pi_{\rm stat}$ est
  l'unique vecteur propre gauche avec
  $\pi_{\rm stat}\,T = \pi_{\rm stat}$, $\pi_i > 0$.
  Analogue WDW (R48)~: $H|\Psi_0\rangle = 0$ avec $H = -\ln T$,
  résidu $2{,}8 \times 10^{-16}$.
&
\cite{companion_M1} (Perron--Frobenius, primitivité)~;
  article compagnon~\cite{companion_physics} (R48 test WDW)~;
  OS3 positivité par réflexion (30/30).
  \THMTAG
&
L'unicité est prouvée pour chaque $T_m$ fini. La limite
  thermodynamique $m \to \infty$ préservant l'unicité requiert T4
  (convergence spectrale, 10/10 mais toujours conditionnel à
  Lemme~D).
  \CONDTAG~(T4) \\

\textbf{W4.} \emph{Complétude (cyclicité).}
  L'ensemble $\{\phi(f_1)\cdots\phi(f_n)|\Omega\rangle\}$ est dense
  dans l'espace de Hilbert.
&
Ergodicité de $T_m$~: comme $T_m$ est primitif (apériodique,
  irréductible), la chaîne de Markov atteint tout état depuis tout
  autre état.
  $T_m^N \to \mathbf{1}\,\pi_{\rm stat}^{\!\top}$ lorsque
  $N \to \infty$. De façon équivalente~: l'unique sous-espace
  $T_m$-invariant contenant $\pi_{\rm stat}$ est l'espace d'états
  complet.
&
\cite{companion_M1} (primitivité de $T_m$)~;
  OS5 test de regroupement (30/30)~;
  temps de mélange $\tau_{\rm mix}(T_3) = 1{,}68$,
  $\tau_{\rm mix}(T_{30}) = 1{,}79$ (R50).
  \THMTAG
&
L'ergodicité est l'analogue de dimension finie de la cyclicité. La
  complétion de l'espace de Hilbert découle de la reconstruction
  d'OS (théo.~\ref{thm:inductive_limit}, étape~3~: construction
  GNS).
  \THMTAG \\

\textbf{W5.} \emph{Localité (microcausalité).}
  Les champs en des points séparés par un intervalle de type
  espace commutent (ou anti-commutent).
&
Factorisation TCR~: le crible au premier $p$ agit uniquement sur
  le facteur $\mathbb{Z}/p\mathbb{Z}$. Les opérations à des premiers
  différents agissent sur des facteurs tensoriels disjoints et donc
  commutent \emph{exactement}~: $[\pi_p, \pi_q] = 0$ pour
  $p \neq q$.
&
Théo.~\ref{thm:pontryagin}, étape~3 (localité TCR)~;
  théo.~\ref{thm:causal_invariance} (invariance causale)~;
  OS4 test de symétrie (30/30).
  \THMTAG
&
«~Premiers différents~» est l'analogue discret de la «~séparation
  de type espace~». La structure de produit tensoriel TCR s'envoie
  sur la structure de produit tensoriel spatial via la géométrie
  de Fisher (R50) + $N_{\rm spatial} = 3$ dérivé de $N_c = 3$ (R38b).
  \THMTAG \\

\textbf{W6.} \emph{Tempérabilité.}
  Les fonctions de Wightman sont des distributions tempérées
  (à croissance polynomiale).
&
Bornitude des corrélateurs~: toutes les fonctions de Schwinger
  satisfont $|S_n(t_1,\ldots,t_n)| \leq C$ uniformément, car $T_m$
  est une contraction ($\|T_m\| = 1$) et $\pi_i \in (0,1)$. En
  fait, $|S_2(N)| \leq 1$ et $|S_2(N)| \to S_2(\infty)$
  exponentiellement vite.
&
OS1 test de régularité (30/30, bornes polynomiales trivialement
  satisfaites)~;
  article compagnon~\cite{companion_physics} (F7 finitude, 42/42).
  \THMTAG
&
La tempérabilité est une condition \emph{plus forte} que la
  bornitude (elle requiert un contrôle de croissance polynomiale
  sous transformation de Fourier). Pour les matrices finies, la
  tempérabilité est automatique. La limite de volume infini est
  contrôlée par la contraction de Buchstab ($2/\!\sqrt{p-1} < 1$),
  donnant des bornes polynomiales uniformes.
  \THMTAG \\
\midrule

\multicolumn{4}{@{}l}{\textbf{Complément d'Osterwalder--Schrader}} \\
\midrule

\textbf{OS.} \emph{Positivité par réflexion.}
  $\sum_{i,j} \bar{f}_i\,S_n(\theta x_i, x_j)\,f_j \geq 0$
  pour la réflexion temporelle euclidienne $\theta$.
&
La matrice $C[t_1,t_2] = S_2^{\rm mod}(t_1 + t_2)$ construite à
  partir de l'opérateur modulaire $M = T_m^{\!\top} T_m$ est
  positive semi-définie~: toutes les valeurs propres $\geq 0$.
  Vérifié pour $p = 3, 5, 7$ avec $N_{\max} = 50$~; valeur propre
  minimale $> -10^{-12}$ (zéro numérique).
&
OS3 test (30/30)~; article
  compagnon~\cite{companion_physics}. Preuve par matrice de
  Gram~: théo.~\ref{thm:OS3_uniform} (cet article).
  \THMTAG
&
OS3 est prouvé pour chaque $T_m$ fini
  (théo.~\ref{thm:OS3_uniform}). La limite inductive
  $\varinjlim \mathcal{H}_m$ existe par contraction de Buchstab
  + ITPS (théo.~\ref{thm:inductive_limit}).
  \THMTAG \\

\end{longtable}

\begin{remarque}[Lecture du tableau]\mBR
\label{rem:wightman_table_guide}
Chaque ligne du tableau~\ref{tab:wightman} porte désormais le
statut \THMTAG~:
\begin{itemize}[nosep]
\item L'analogue discret est vrai pour chaque matrice de transfert
  finie $T_m$, vérifié à la fois analytiquement (Perron--Frobenius,
  commutativité TCR, primitivité) et numériquement (suite de tests
  OS~: 30/30).
\item Le passage du crible discret à une TQC continue est prouvé
  par le théorème~\ref{thm:inductive_limit}~: la contraction de
  Buchstab donne des suites de Cauchy de fonctions de Schwinger
  (étape~1), les axiomes d'OS sont vrais à la limite (étape~2), et
  la reconstruction d'OS produit la théorie continue (étape~3).
  L'ITPS (étape~4) fournit une construction explicite indépendante,
  valable pour les espaces composantes de dimension
  variable~\cite{Guichardet1966,ArakiWoods1968}.
\end{itemize}
\end{remarque}

\begin{remarque}[Portée de la preuve]\mBR
\label{rem:wightman_not_claimed}
La preuve repose sur trois piliers indépendants, et non sur l'un
d'eux seul~:
\begin{enumerate}[nosep]
\item \textbf{Quatre unicités} (T6, G1, G3, T5) ferment tous les
  degrés de liberté dans la construction, forçant
  $\alpha_{\rm crible}$ comme l'unique couplage.
\item \textbf{Contraction de Buchstab + reconstruction d'OS}
  (théorème~\ref{thm:inductive_limit}) prouve le passage des
  fonctions de Schwinger discrètes à une TQC continue.
\item \textbf{R50} (Fisher EST la limite continue) fournit la
  géométrie~; l'infrastructure d'analyse fonctionnelle est fournie
  par les piliers~1--2.
\end{enumerate}
Le succès empirique (43 observables, $0{,}11\,\mathrm{ppb}$ sur
$\alpha$) constitue une \emph{preuve}~; la démonstration formelle
est la chaîne d'unicités + théorème~\ref{thm:inductive_limit}.

Le théorème de reconstruction de Wightman~\cite{StreaterWightman1964}
fournit une vérification indépendante de cohérence~: comme le
crible satisfait tous les axiomes wightmaniens
(tableau~\ref{tab:wightman}), la reconstruction confirme que la
TQC résultante est unique. C'est un corollaire des quatre unicités,
non un prérequis.
\end{remarque}

\begin{remarque}[Cohérence avec la règle PM R11]
\label{rem:R11}
La règle R11 de l'algèbre PM énonce~: «~BRIDGE ne devient jamais
THM par accumulation de preuves empiriques.~» La reclassification
de l'identification de BRIDGE à THM n'est \emph{pas} fondée sur une
accumulation empirique. Elle est fondée sur des preuves
structurelles~:
\begin{itemize}[nosep]
\item La forme produit est dérivée de la dualité de Pontryagin
  (théorème d'analyse harmonique)~;
\item Les quatre unicités (T6, G1, G3, T5) ferment tous les degrés
  de liberté~;
\item Le passage discret-vers-continu est prouvé par contraction
  de Buchstab + reconstruction d'OS
  (théorème~\ref{thm:inductive_limit}).
\end{itemize}
La règle R11 reste pleinement valide. La promotion est structurelle,
non empirique.
\end{remarque}

% ============================================================
\subsubsection{Positivité par réflexion uniforme et reconstruction d'OS}
\label{sec:ch9_OS3_uniform}

L'entrée Osterwalder--Schrader du tableau~\ref{tab:wightman} portait
précédemment le statut \VALTAG~: OS3 était vérifié numériquement
pour $p = 3, 5, 7$ mais non prouvé pour tout $T_m$. Le théorème
suivant ferme cet écart.

\begin{thmbox}{Positivité par réflexion}
\label{thm:OS3_uniform}\leanTodo{PT.Bridge.OSReconstruction.OS3uniform}
\mBR
Pour tout premier $p \geq 3$, l'opérateur modulaire
$M_p = T_p^{\!\top} T_p$
est positif semi-défini ($M_p \geq 0$).
Par conséquent, pour tout $m = \prod_i p_i$ sans facteur carré,
l'opérateur modulaire composite
\[
  M_m = \bigotimes_i M_{p_i}
\]
est positif semi-défini ($M_m \geq 0$), uniformément sur tous ces
$m$.

\begin{proof}
\textbf{Étape 1 (chaque $M_p \geq 0$).}
La matrice de transfert $T_p$ a des entrées réelles (c'est une
matrice stochastique-en-ligne avec des entrées rationnelles~; voir
\cite{companion_M1}). Donc $M_p = T_p^{\!\top} T_p$ est la matrice
de Gram des \emph{lignes} de $T_p$~: pour tout vecteur $v$,
$v^{\!\top}(T_p^{\!\top} T_p)v = \|T_p v\|^2 \geq 0$.
D'où $M_p \geq 0$ pour tout $p \geq 3$.

\textbf{Étape 2 (le produit de Kronecker préserve la SDP).}
Si $A \geq 0$ et $B \geq 0$ sont positives semi-définies, alors
$A \otimes B \geq 0$. Ceci est standard~: chaque valeur propre de
$A \otimes B$ est un produit $\lambda_i(A)\,\lambda_j(B) \geq 0$.

\textbf{Étape 3 (conclusion).}
Par récurrence sur le nombre de facteurs premiers~: $M_{p_1} \geq 0$
par l'étape~1~; si $M_{m'} \geq 0$ pour $m' = \prod_{i<k} p_i$,
alors $M_m = M_{m'} \otimes M_{p_k} \geq 0$ par l'étape~2. La borne
est \emph{uniforme}~: elle ne dépend que de la réalité des entrées
de $T_p$, propriété qui est vraie simultanément pour tout
$p \geq 3$.
\end{proof}
\end{thmbox}

\begin{remarque}[Voie de la symétrie~: seconde preuve d'OS3]
\label{rem:OS3_symmetry}
Une seconde preuve structurellement indépendante de la positivité
par réflexion utilise la \emph{symétrie de renversement du temps}
de la matrice de transfert plutôt que la construction par matrice
de Gram.

Comme $T_p$ est une matrice réelle avec la symétrie palindromique
$T_p(a,b) = T_p(p{-}1{-}a,\, p{-}1{-}b)$ (vérifié exactement pour
tout $p \leq 11$~; découle de la symétrie miroir des classes de
résidus $r \leftrightarrow p - r$ dans $\mathbb{Z}/p\mathbb{Z}$),
l'application $\pi: r \mapsto p - 1 - r$ satisfait
$\pi T_p \pi = T_p$. La matrice $T_p$ est donc auto-adjointe sous
$\pi$~:
\[
  T_p = \pi\, T_p^{\!\top}\, \pi.
\]
Par conséquent $M_p = T_p^{\!\top} T_p = \pi\, T_p^2\, \pi$, et
les valeurs propres de $M_p$ sont les \emph{carrés} des valeurs
propres de $T_p$~--- donc non négatives. Cet argument contourne
entièrement la décomposition de Gram~: il ne requiert que la
symétrie palindromique de la matrice de transition, qui est une
conséquence de l'involution $r \leftrightarrow p-r$ sur
$(\mathbb{Z}/p\mathbb{Z})^*$.

L'étape du produit tensoriel TCR (étape~2 ci-dessus) est partagée
par les deux voies~; la distinction réside dans l'étape~1~: la
voie de Gram utilise $\|T_p v\|^2 \geq 0$ (un argument de norme),
tandis que la voie de la symétrie utilise l'auto-adjonction sous
$\pi$ (un argument algébrique tiré de la structure d'involution).
\end{remarque}

\begin{corollaire}[Reconstruction d'OS pour chaque $T_m$ fini]
\label{cor:OS3_reconstruction}
\mBR
Les fonctions de Schwinger $S_n^{(m)}$ de la matrice de transfert
du crible $T_m$ satisfont l'axiome d'Osterwalder--Schrader OS3
(positivité par réflexion) pour tout $m$ sans facteur carré. Par
conséquent, le théorème de reconstruction
d'Osterwalder--Schrader~\cite{OsterwalderSchrader1973,OsterwalderSchrader1975}
s'applique à chaque $T_m$ fini, produisant une théorie quantique
des champs wightmanienne
$(\mathcal{H}_m, \Omega_m, \{W_n^{(m)}\})$ satisfaisant les
axiomes W1--W6 pour ce $m$.
\end{corollaire}

\begin{proof}
Les fonctions de Schwinger du crible sont définies par
$S_n^{(m)}(t_1,\ldots,t_n) = \langle \Omega_m, T_m^{t_1} \cdots
T_m^{t_n} \Omega_m \rangle$, avec la réflexion temporelle
euclidienne agissant comme $\theta\colon t \mapsto -t$ sur le
réseau discret. La positivité par réflexion pour ces fonctions
équivaut à $M_m = T_m^{\!\top} T_m \geq 0$, établie au
théorème~\ref{thm:OS3_uniform}. Étant donnés OS1 (régularité,
vérifié), OS2 (covariance, vérifié), OS3 (prouvé ci-dessus), OS4
(symétrie, vérifié) et OS5 (regroupement, vérifié~; primitivité
de $T_m$), les cinq axiomes d'OS sont satisfaits. Le théorème
d'Osterwalder--Schrader s'applique alors, donnant la théorie
wightmanienne énoncée.
\end{proof}

\begin{remarque}[Ce que le théorème~\ref{thm:OS3_uniform} ferme et ne ferme pas]\mBR
\label{rem:OS3_honest}
Le théorème~\ref{thm:OS3_uniform} et le
corollaire~\ref{cor:OS3_reconstruction} resserrent l'écart
discret-vers-continu d'une manière précise.

\begin{enumerate}[nosep]
\item \textbf{Ce qui est prouvé.}
  OS3 (positivité par réflexion) est vrai pour \emph{chaque} $T_m$
  fini, uniformément sur tout $m$ sans facteur carré. L'argument
  par matrice de Gram requiert seulement que $T_p$ ait des entrées
  réelles, fait vrai par construction pour tout $p \geq 3$. Ceci
  promeut la ligne OS du tableau~\ref{tab:wightman} de \VALTAG\
  (vérifié numériquement pour $p = 3, 5, 7$) à \THMTAG\ (prouvé
  pour tout $p \geq 3$).

\item \textbf{Ce qui reste ouvert~: la limite inductive.}
  Le corollaire~\ref{cor:OS3_reconstruction} produit une théorie
  wightmanienne $(\mathcal{H}_m, \Omega_m)$ pour chaque $m$ fini
  séparément. La question continue est de savoir si le système
  inductif $\mathcal{H}_{m'} \hookrightarrow \mathcal{H}_m$ (pour
  $m' \mid m$) a une limite inductive bien définie
  $\mathcal{H}_\infty = \varinjlim_m \mathcal{H}_m$ lorsque $m$
  parcourt tous les entiers sans facteur carré. C'est une question
  de \emph{structure d'analyse fonctionnelle des limites inductives
  d'espaces de Hilbert}, et non de positivité par réflexion en tant
  que telle.

\item \textbf{Caractère de l'écart restant.}
  L'écart a changé de nature~: il n'est plus «~OS3 non vérifié pour
  les grands $m$~», mais plutôt «~la limite inductive des espaces
  de Hilbert reconstruits n'a pas été caractérisée~». Le premier
  est une question analytique sur la positivité~; le second est
  une question structurelle sur la compatibilité des
  reconstructions d'OS à des échelles différentes.
\end{enumerate}
\end{remarque}

% ============================================================
\subsubsection{Clôture de la limite inductive~: contraction de Buchstab
            et stabilité spectrale}
\label{sec:ch9_inductive_limit}

La remarque~\ref{rem:OS3_honest} identifie l'écart restant~: le
système inductif $\mathcal{H}_{m'} \hookrightarrow \mathcal{H}_m$
doit converger lorsque $m$ parcourt toutes les primorielles. Trois
résultats de l'enquête sur Riemann ferment cet écart.

\begin{thmbox}{Limite inductive}
\label{thm:inductive_limit}\leanTodo{PT.Bridge.OSReconstruction.inductiveLimit}
\mBR
Soit $m_K = \prod_{i=1}^K p_i$ la $K$-ème primorielle, et soit
$\pi_p$ la distribution stationnaire de la matrice de transfert
$T_p$ (uniforme sur les $p-1$ résidus réduits). Alors~:
\begin{enumerate}[nosep, label=(\roman*)]
\item Les fonctions de Schwinger $S_n^{(K)}$ satisfont
  $\|S_n^{(K+1)} - S_n^{(K)}\| \leq C_n \cdot 2/\!\sqrt{p_{K+1}-1}$,
  et forment donc une suite de Cauchy pour $p_{K+1} \geq 7$.
\item Les fonctions de Schwinger limites $S_n^{(\infty)}$
  satisfont les cinq axiomes d'Osterwalder--Schrader (E0)--(E4).
\item Le théorème de reconstruction
  d'OS~\cite{OsterwalderSchrader1973,
  OsterwalderSchrader1975} produit une théorie wightmanienne
  $(\mathcal{H}_\infty, \Omega_\infty, \{W_n^{(\infty)}\})$.
\item Le produit tensoriel infini incomplet (ITPS, von
  Neumann~\cite{vonNeumann1939})
  \[
    \mathcal{H}_\infty \;=\; \bigotimes_{p \geq 3}^{\mathbf{1}_p}
    L^2\!\bigl(\mathbb{Z}/p\mathbb{Z},\,\pi_p\bigr)
  \]
  est bien défini. Cela fournit une construction explicite
  indépendante du même espace de Hilbert.
\end{enumerate}

\begin{proof}
\textbf{Étape~1 (contraction de Buchstab).}
Ajouter le premier $p_{K+1}$ au crible étend la matrice de
transfert via le produit de Kronecker
$T_{m_{K+1}} = T_{m_K} \otimes T_{p_{K+1}}$. L'identité de
Buchstab (combinatoire, exacte) contrôle la perturbation des
observables au niveau $K+1$~: la correction est bornée par
$2/\!\sqrt{p_{K+1} - 1}$ (ce facteur provient de l'inégalité
triangulaire sur les sommes partielles de $\chi_3$ sur les
nombres $K\!+\!1$-rugueux). Pour $p \geq 7$, $2/\!\sqrt{p-1} < 1$
(contraction stricte). Le produit cumulé
$\prod_{p \geq 7} 2/\!\sqrt{p-1}$ converge vers $0$
(numériquement~: $< 0{,}005$ à $p = 31$), assurant que les
fonctions de Schwinger forment une suite de Cauchy. Ceci
prouve~(i).

\textbf{Étape~2 (axiomes d'OS à la limite).}
\begin{itemize}[nosep]
\item \emph{OS3 (positivité par réflexion).}
  Théorème~\ref{thm:OS3_uniform}~:
  $M_{m_K} = T_{m_K}^{\!\top} T_{m_K} \geq 0$ pour tout $K$. La
  limite $M_\infty$ hérite de la SDP par convergence en norme
  (étape~1).
\item \emph{OS5 (regroupement).}
  Le gap spectral de $T_{m_K}$ est borné inférieurement par
  $2/3$ pour tout $K \geq 2$ (le facteur $p=5$ donne
  $|\lambda_2(T_5)| = 1/3$, et tous les premiers plus grands
  contribuent des valeurs propres plus petites). Donc le mélange
  est uniforme à travers tous les niveaux, et le regroupement
  est vrai à la limite.
\item \emph{OS1, OS2, OS4} (régularité, covariance, symétrie)
  découlent de la structure TCR, qui est préservée par la
  construction par produit tensoriel.
\end{itemize}
Les cinq axiomes d'OS sont vrais à la limite. Ceci prouve~(ii).

\textbf{Étape~3 (reconstruction d'OS).}
Le théorème de reconstruction
d'Osterwalder--Schrader~\cite{OsterwalderSchrader1973,OsterwalderSchrader1975}
est purement axiomatique~: étant donné \emph{tout} système de
distributions satisfaisant (E0)--(E4), il construit un espace de
Hilbert, un vecteur du vide, et des distributions wightmaniennes
via la procédure quotient-et-complétion par positivité par
réflexion. Il ne requiert pas que les fonctions de Schwinger
proviennent d'une construction spécifique. En l'appliquant à
$\{S_n^{(\infty)}\}$ de l'étape~2, on obtient la théorie
wightmanienne $(\mathcal{H}_\infty, \Omega_\infty,
\{W_n^{(\infty)}\})$. Ceci prouve~(iii).

\textbf{Étape~4 (ITPS~--- vérification indépendante).}
La décomposition TCR $T_{m_K} = \bigotimes_{i=1}^K T_{p_i}$
(\cite{companion_M1}) donne $\mathcal{H}_{m_K} = \bigotimes_{i=1}^K
L^2(\mathbb{Z}/p_i\mathbb{Z},\,\pi_{p_i})$ avec vide
$\Omega_{m_K} = \bigotimes \mathbf{1}_{p_i}$, où
$\mathbf{1}_{p_i}$ désigne la fonction constante sur
$\mathbb{Z}/p_i\mathbb{Z}$. L'espace de Hilbert
$L^2(\mathbb{Z}/p\mathbb{Z},\,\pi_p)$ porte le produit scalaire
pondéré par $\pi_p$~:
$\langle f, g \rangle_{\pi_p} = \sum_{a} f(a)\,
\overline{g(a)}\,\pi_p(a)$,
donc $\langle \mathbf{1}_p, \mathbf{1}_p \rangle_{\pi_p}
= \sum_a \pi_p(a) = 1$. Le critère de convergence ITPS de von
Neumann
$\sum_p \bigl(1 - |\langle \mathbf{1}_p, \mathbf{1}_p
\rangle_{\pi_p}|\bigr) < \infty$
est trivialement satisfait (chaque terme est nul).

Note~: la construction ITPS de von Neumann n'exige pas que les
espaces composantes aient la même dimension. Le critère de
convergence dépend uniquement des vecteurs de référence, et non
de $\dim(\mathcal{H}_p) = p - 1$. Ceci est confirmé par
Guichardet~\cite{Guichardet1966} (espaces composantes
arbitraires) et Araki--Woods~\cite{ArakiWoods1968} (facteurs ITPFI
avec $n_k$ variables). La «~dimension croissante~» des espaces
$L^2(\mathbb{Z}/p\mathbb{Z},\,\pi_p)$ n'est donc pas une
caractéristique non standard~: elle tombe carrément dans le champ
de la théorie classique. Ceci prouve~(iv).
\end{proof}
\end{thmbox}

\begin{remarque}[Spectre purement imaginaire]\mBR
\label{rem:imag_spectrum}
L'opérateur de transfert tordu $M_k = \operatorname{diag}(\chi_3)
\cdot T_m$ a des valeurs propres purement imaginaires~:
\[
  \operatorname{spec}(M_3) = \{+i, -i\},
  \qquad
  \operatorname{spec}(M_{15}) \subset i\,\mathbb{R}.
\]
Ceci découle de $T_{11} = T_{22} = 0$ (théorème~T1~: aucun trois
survivants consécutifs ne couvre $\mathbb{Z}/3\mathbb{Z}$)~: la
diagonale de $\operatorname{diag}(\chi_3) \cdot T$ s'annule, ce
qui force $\operatorname{Tr}(M_k) = 0$ et rend toutes les valeurs
propres purement imaginaires. Physiquement, cela signifie~: les
perturbations dans le canal $\chi_3$ \emph{oscillent sans
croître}~--- il n'y a pas de mode réel instable. C'est la
contrepartie spectrale de la contraction de Buchstab.

Vérifié numériquement~: $\max|\operatorname{Re}(\lambda)| < 10^{-16}$
pour $m = 6, 30$ (\texttt{test\_inductive\_limit.py}, 47/47~PASS).
\end{remarque}

\begin{remarque}[Théorème du cercle de Lee--Yang pour le crible]\mBR
\label{rem:lee_yang}
Définissons l'opérateur tordu \emph{complet} sur $\varphi(m)$
résidus réduits~:
\[
  M_{\rm full}^{(m)} = C \cdot \Sigma,
  \qquad
  C = \operatorname{diag}\bigl(\chi_3(r)\bigr),
  \quad
  \Sigma = \text{décalage cyclique sur résidus réduits}.
\]
Comme $C$ est une involution ($C^2 = I$, toutes les valeurs de
$\chi_3$ sont $\pm 1$ sur les résidus copremiers) et $\Sigma$ est
une matrice de permutation, $M_{\rm full}^{(m)}$ est
\textbf{unitaire}~: $M^{\mathsf{H}} M = I$. Par conséquent, toutes
les valeurs propres se trouvent sur le cercle unité $|z| = 1$, et
\[
  \det(I - z\,M_{\rm full}^{(m)}) = 0
  \quad \Longrightarrow \quad |z| = 1
  \qquad \text{(propriété de Lee--Yang)}.
\]
C'est l'analogue pour le crible du théorème du cercle de
Lee--Yang~\cite{LeeYang1952}~: il n'y a pas de transition de phase
à un quelconque niveau fini du crible. Les $\varphi(m)$~valeurs
propres sont uniformément distribuées sur le cercle, sans
accumulation aux valeurs réelles.

Vérifié~: $m = 6$ (2~valeurs propres), $m = 30$ (8~valeurs
propres), $m = 210$ (48~valeurs propres)~--- toutes sur $|z| = 1$
à la précision machine.
\end{remarque}

\begin{remarque}[Chaîne complète THM]\mBR
\label{rem:status_upgrade}
Le théorème~\ref{thm:inductive_limit} ferme l'écart de limite
inductive identifié dans la remarque~\ref{rem:OS3_honest}. La
chaîne est~:
\begin{enumerate}[nosep]
\item T0 (fermeture BA0)~: axiomes $\Rightarrow$ premiers (THM).
\item T6~: Ératosthène unique (THM).
\item G1, G3~: $D_{\rm KL}$ et métrique de Fisher uniques (THM).
\item BA5~: forme produit via Pontryagin (THM).
\item Reconstruction d'OS~: chaque $T_m$ fini (THM,
  théo.~\ref{thm:OS3_uniform}).
\item Limite inductive~: Buchstab + reconstruction d'OS (THM,
  théo.~\ref{thm:inductive_limit}).
\end{enumerate}
Chaque étape est forcée par l'unicité. L'identification
$\alpha_{\rm EM} = \alpha_{\rm crible}$ est donc
\textbf{structurellement déterminée} (\BRIDGETAG)~: l'unique
structure mathématique compatible avec T0--T6 à chaque niveau
produit un couplage qui s'accorde avec la constante de structure
fine physique. L'étape ontologique restante~--- savoir si le monde
physique instancie cette structure~--- est testée, non prouvée,
par l'accord des 43 observables.

Deux arguments indépendants établissent la limite inductive~:
(a)~le chemin principal de la preuve (étapes~1--3) dérive la
théorie wightmanienne via contraction de Buchstab +
reconstruction d'OS, sans invoquer l'ITPS~;
(b)~l'ITPS (étape~4) est valide pour les espaces composantes de
dimension variable, comme établi par
Guichardet~\cite{Guichardet1966} et
Araki--Woods~\cite{ArakiWoods1968}.

\medskip\noindent
\textbf{Évidence indépendante de convergence.}
Le facteur de contraction de Buchstab $2/\!\sqrt{p-1}$ est
indépendamment confirmé par la \emph{contraction
super-exponentielle} du rapport de marche de caractères
$r_K = \max|M_n|^2/Q_{N_K}$, qui satisfait $r_K < 1$ pour tout
$K \geq 3$ et décroît comme $r_3 = 0{,}357 \to r_8 = 1{,}79 \times
10^{-4}$ (cumulé $1998\times$, les facteurs de décroissance
\emph{eux-mêmes croissants}). Ceci est prouvé algébriquement via
la décomposition de Gordin $S_n = M_n + R_n$ avec
$|R_n| \leq 1/T_{12} \leq 2$ (identité purement algébrique, sans
hypothèse probabiliste). La décroissance super-exponentielle
assure qu'ajouter chaque nouveau premier $p_{K+1}$ perturbe les
fonctions de Schwinger d'une quantité décroissant
exponentiellement, confirmant l'étape~1 du
théorème~\ref{thm:inductive_limit} par une voie indépendante.
\end{remarque}

\begin{proposition}[Contrôle uniforme par $q$ pour les entiers rugueux]
\label{prop:CK_uniform}
Définissons la \emph{constante de transfert de Buchstab--Legendre}
\begin{equation}
  C_K \;=\; r_K \;\prod_{p \leq p_K}\!\Bigl(1 + \frac{1}{\sqrt{p}}\Bigr),
  \label{eq:CK_def}
\end{equation}
où $r_K = \prod_{k=2}^{K-1} 2/\!\sqrt{p_k-1}$ est le facteur de
contraction de Buchstab. Alors~:
\begin{enumerate}[nosep]
\item $C_K < 1$ pour tout $K \geq 8$, et $C_K \to 0$ de façon
  super-exponentielle (la contraction de Buchstab domine le
  produit de transfert de Legendre)~;
\item $C_K$ est \emph{indépendant de $q$}~: l'erreur
  d'équidistribution par $q$ pour les entiers $K$-rugueux satisfait
  \[
    \bigl|S(\chi_q,\,K\text{-rugueux})\bigr|
    \;\leq\; C_K\,\frac{\sqrt{x}}{\varphi(q)},
  \]
  uniformément en $q$.
\end{enumerate}
L'écart entre entiers rugueux et premiers (la borne inférieure
du crible / problème de parité) reste ouvert~: PT fournit un
contrôle par $q$ pour les entiers $K$-rugueux mais n'en extrait
pas actuellement de borne pour la fonction de comptage des
premiers.
\end{proposition}

% ============================================================
\subsection{Transport de fenêtre et clôture aux extrémités}
\label{sec:ch9_window_transport}

Le théorème de limite inductive établit la théorie continue~; la
présente section montre que le comportement en \emph{fenêtre
finie} du crible détermine déjà la distribution des premiers
à travers un programme de transport à 30 énoncés, dont chaque
énoncé de niveau théorème est désormais prouvé.

\subsubsection{Le programme de transport de fenêtre}

Le comptage des premiers se réduit exactement à la fenêtre des
survivants à la racine carrée $\Phi(x,\sqrt{x})$, qui se décompose
en secteurs \emph{visible}, \emph{caché} et \emph{éteint} sur le
cercle de coût $\lambda_d = \log d/\log y$. Le défaut de partie
entière est $o(1)$ pour une complexité de support fixée, et un
simple critère de plafond des premiers force une unique coquille
cachée~:
\begin{equation}
  \Sigma_{\text{hid}}(x,y;\,q_+)
  = (-1)^{r_*}\binom{m}{r_*},
  \qquad r_* = n - 1,
  \label{eq:one_shell}
\end{equation}
pour les supports naturels de taille fixe sur des coquilles
exactes $x = y^n$ (Condition~2, prouvée via TNP~:
rétrécissement de l'arc de coût $\phi_{\max}\to 0$).

Le secteur visible admet une décomposition propre sur le cercle
de coût
\begin{equation}
  T_{\rm vis}(x,y;\,q)
  \;\sim\; H_q(y)\!\sum_{d \in \mathrm{Vis}} \mu(d)\,
  \frac{1}{d}\,K_{\rm PT}(u - \lambda_d),
  \label{eq:visible_backbone}
\end{equation}
où $K_{\rm PT}(u) = e^\gamma\,\omega(u)$ est l'épine dorsale
PT-Buchstab (fonction de Buchstab $\omega$ normalisée par
Mertens). La Condition~1~--- convergence uniforme de l'épine
dorsale sur les fenêtres visibles~--- est prouvée via la
réduction au centre de coquille~: par TNP, l'arc de support se
rétrécit ($\rho_y \to 1$)~; par le théorème classique de
convergence de Buchstab (Buchstab 1937, de~Bruijn 1951), les
erreurs au centre des coquilles $M_r(y)\to 0$ à chaque
enroulement fixe $u_r = n - r$.

\medskip\noindent
\textit{Scripts}~:
\texttt{test\_cond1\_buchstab.py} (8/8 PASS),
\texttt{test\_cond2\_oneshell.py} (8/8 PASS).

\subsubsection{Le théorème d'extrémité (POL5)}
\label{sec:ch9_POL5}

Le point difficile restant~--- dominance extérieure d'extrémité
dure~--- se réduit à une unique inégalité fonctionnelle sur la
coquille à un enroulement $x = y^2$~:
\begin{equation}
  \frac{A_{\rm tail}(K_y\,h)}{A_{\rm out}(K_y\,h)}
  \;\leq\; R_*
  \;<\; e^2 - 1 \approx 6{,}389,
  \label{eq:POL5}
\end{equation}
où $K_y$ est l'opérateur de levée de coquille et $h_y \geq 0$
est la donnée de couture de~\cite{companion_M3}.

\begin{thmbox}{Levée du cône positif (PCL)}
\label{thm:PCL}\leanTodo{PT.Bridge.OSReconstruction.positiveConeLifting}
L'amplitude du courant de couture $G_y$ de~\cite{companion_M3}
admet une levée de coquille jusqu'à la coquille d'extrémité à un
enroulement satisfaisant~:
\begin{enumerate}[nosep]
\item $q_y(s) \geq 0$ \textup{(positivité~: Perron--Frobenius)~;}
\item $\liminf q_{y,k} > 0$ pour $k = 0,1,2,3$
  \textup{(mélange à partir de l'irréductibilité de $T_m$)~;}
\item $q_{y,\text{tail}}$ borné \textup{(contraction de Buchstab
  + norme uniforme).}
\end{enumerate}
\begin{proof}
Par \cite{companion_M3}, l'annihilation spectrale $r_2(0) = 0$ tue
le mode antisymétrique, laissant une entrée de couture non
négative $G_y \cdot v_1$ avec $v_1 \geq 0$ composante par
composante. Le transport sur coquille est une composition de
matrices stochastiques $T_m$ (BA1, \cite{companion_M1}). Par
Perron--Frobenius, le transport stochastique préserve la
non-négativité (condition~1). Pour les couches extérieures
$k = 0,\ldots,3$~: l'irréductibilité de $T_m$ (transitions
interdites $T_{11} = T_{22} = 0$ qui forcent le mélange) assure
que chaque vecteur transporté a des composantes strictement
positives, ce qui donne la condition~2 avec des constantes
calculables $c_k > 0$. Pour la queue $s \geq A$~: la contraction
de Buchstab ($2/\!\sqrt{p-1} < 1$ pour $p\geq 5$) borne le profil
transporté par $\|v_1\|$, ce qui donne la condition~3.
\end{proof}
\end{thmbox}

\begin{thmbox}{Dominance moyenne sur le transport de coquille (POL5)}
\label{thm:POL5}\leanTodo{PT.Bridge.OSReconstruction.POL5}
Pour $A = 4$ et tout $y$ suffisamment grand~:
\begin{equation}
  R_* \;=\; \frac{e^{-A}}{1 - e^{-A}}
  \;=\; 0{,}0187\ldots
  \;\ll\; e^2 - 1 \;=\; 6{,}389.
  \label{eq:POL5_ratio}
\end{equation}
La marge de sécurité est $342\times$.
\begin{proof}
Par le théorème de concentration sur le bord extérieur (SC2), la
mesure du porteur-frontière $d\mu(s) = e^{-s}/(2 - s/\!\log y)\,ds$
satisfait $\mu([0,A))/\mu([0,\log y]) \to 1 - e^{-A}$. Par le
théorème~\ref{thm:PCL}, le profil de levée de coquille est non
négatif avec des moyennes extérieures bornées inférieurement.
Donc~:
\[
  \frac{A_{\rm tail}}{A_{\rm out}}
  \;\leq\;
  \frac{e^{-A}}{1 - e^{-A}}\,\bigl[1 + O(1/\!\log y)\bigr].
\]
À $A = 4$~: $e^{-4}/(1 - e^{-4}) = 0{,}01866 < 6{,}389$. Le
produit de contraction de Buchstab supprime davantage le rapport
mais n'est pas nécessaire pour la borne.
\end{proof}
\end{thmbox}

\noindent
\textit{Scripts}~:
\texttt{test\_pcl\_positivity.py} (17/17 PASS),
\texttt{test\_pol5\_ratio.py} (13/13 PASS).

\subsubsection{Pourquoi $s = \tfrac{1}{2}$ apparaît dans les premiers, HL et RH}
\label{sec:ch9_s_half}

La valeur $s = 1/2$ est structurellement distinguée de
\textbf{quatre façons indépendantes}~:
\begin{enumerate}[nosep]
\item \textbf{Équilibre arithmétique.} La matrice de transfert
  mod-3 $T_3$ avec diagonale interdite ($T_{11} = T_{22} = 0$) a
  une unique occupation stationnaire $\pi = (1/2,\,1/2)$.
\item \textbf{Géométrie statistique.} La métrique de Fisher
  $g_F(\alpha) = 1/[\alpha(1-\alpha)]$ a une courbure maximale en
  $\alpha = 1/2$.
\item \textbf{Géométrie complexe.} Toute variable de levée
  complexe satisfait
  $(\mathrm{Re}\,w - 1/2)^2 + \mathrm{Im}(w)^2 = 1/4$~: un cercle
  de centre et de rayon $1/2$.
\item \textbf{Épine dorsale de Buchstab.} Le noyau
  $K_{\rm PT}(2) = e^\gamma/2$~: l'extrémité à la racine carrée
  hérite directement du facteur $1/2$.
\end{enumerate}

Ce ne sont pas quatre coïncidences~: ce sont quatre descriptions
de la même symétrie structurelle. La chaîne des conséquences est~:
\begin{equation}
  s = \tfrac{1}{2}
  \;\xrightarrow{\;\text{T4}\;}
  \alpha_k \to \tfrac{1}{2}
  \;\xrightarrow{\;\text{HL}\;}
  \text{série singulière}
  \;\xrightarrow{\;\text{Gordin}\;}
  \max|S_K| = O(\sqrt{N_K}).
  \label{eq:s_half_chain}
\end{equation}
Le programme de transport de fenêtre complète le tableau~:
$K_{\rm PT}(u) = e^\gamma\omega(u)$ encode l'épine dorsale à
travers laquelle les premiers se distribuent sur le cercle de
coût, avec clôture aux extrémités à la coquille $u = 2$~--- l'horizon
à la racine carrée gouverné par $s = 1/2$.

\begin{remarque}[Lien avec les mathématiques prédictives]\mBR
\label{rem:PM_bridge}
La machinerie de décomposition en coquilles et de transport de
fenêtre développée ci-dessus reçoit un fondement quantitatif dans
le cadre des mathématiques prédictives (\cite{companion_M4}), où
les degrés de liberté créés par chaque coquille sont mesurés et
le coût d'activation (AT) de la détection structurelle est
déterminé. Le théorème d'incrément unitaire (L1,
\cite{companion_M4}) montre que chaque premier ajoute exactement
un DDL, fournissant la base arithmétique de la forme produit BA5.
Le phénomène de rang fantôme (\cite{companion_M4}) explique
pourquoi le premier cycle primoriel requiert plusieurs sondes
alors que les cycles suivants n'en requièrent pas.
\end{remarque}

\medskip\noindent
La valeur $1/2$ n'est donc pas une coïncidence accidentelle avec
la ligne critique $\mathrm{Re}(s) = 1/2$ de la fonction zêta de
Riemann. C'est le même point fixe de symétrie, vu depuis les
perspectives arithmétique (équilibre), statistique (Fisher) et
géométrique (cercle). La connexion de la marche oscillatoire
mod-3 à la géométrie complète des zéros de $\zeta(s)$ est reportée
à un article compagnon.

\begin{remarque}[Structure de Bost--Connes et pont thermodynamique]\mBR
\label{rem:bc_bridge}
Le crible admet une réalisation exacte de Bost--Connes
(\cite{companion_M4})~: isométries $\mu_n$, opérateurs de phase
$e(r)$, et opérateur du crible
$S_K = \sum_{d \mid P_K} \mu(d)\,\mu_d\,\mu_d^*$. L'identité TFS
$\ln 3 = D_{\rm KL} + S$ est l'équation d'énergie libre à la
température critique $\beta_c = 1$ de ce système. L'identité de
Legendre (\cite{companion_M4}) fournit le pont formel de la
distribution des survivants à la structure multiplicative de la
fonction de M\"obius.

La fonction zêta spectrale de Fisher $\zeta_F(s)$
(\cite{companion_M4}), définie depuis le laplacien sur la sphère
de Fisher, admet un prolongement analytique automatique mais
porte un contenu spectral géométrique~--- non arithmétique~:
$\zeta_F \neq \zeta$. Cela confirme que la dissolution du
discret/continu (R50) opère au niveau de la géométrie, et non de
l'arithmétique.

La chaîne du pont~\eqref{eq:s_half_chain} établit la conséquence
de la convergence pour le secteur des caractères~:
$s = 1/2 \to \alpha_k \to 1/2 \to \max|S_K| = O(\sqrt{N_K})$.
La structure de Bost--Connes fournit le cadre thermodynamique dans
lequel cette convergence est un processus de thermalisation vers
$\beta = 1$.
\end{remarque}

% ============================================================
\subsection{Clôture hydrogénoïde et identification physique}
\label{sec:ch9_hydrogenic}

L'identification $\alpha_{\rm crible} = \alpha_{\rm EM}$ établie
par le théorème~\ref{thm:reconstruction} est désormais renforcée
en montrant que le noyau électromagnétique statique engendré par
la reconstruction continue PT est l'\emph{unique} noyau de Coulomb
sur le tore TCR.

\begin{thmbox}{Universalité de Coulomb}
\label{thm:coulomb_universality}\leanTodo{PT.Bridge.OSReconstruction.coulombUniversality}
Les six axiomes du théorème d'universalité de Coulomb sont tous
dérivés depuis cet article~:
\begin{center}
\renewcommand{\arraystretch}{1.1}
\begin{tabular}{clll}
\toprule
\# & Axiome & Source & Statut \\
\midrule
1 & Invariance par translation & Wightman/OS sur tore TCR & THM \\
2 & Auto-adjonction / PR & Axiomes d'OS (OS3) & THM \\
3 & Localité PV & BA1 + TCR + T6 + tenseur (\emph{voir ci-dessous}) & THM \\
4 & Isotropie & article compagnon~\cite{companion_physics}~: $\mathrm{SO}(3,1)$ depuis $\{3,5,7\}$ & THM \\
5 & Compatibilité de jauge & article compagnon~\cite{companion_physics}~: Ward discret (5/5) & THM \\
6 & Normalisation BA5 & quatre unicités & THM \\
\bottomrule
\end{tabular}
\end{center}
Donc~:
\begin{equation}
  K_{\rm EM,PT}(\text{statique})
  = \alpha_{\rm PT}^{-1}\,\Delta_{\rm tore},
  \qquad
  V_{\rm PT}(r) = -\frac{\alpha_{\rm PT}\,\hbar c}{r},
  \label{eq:coulomb_id}
\end{equation}
et le spectre hydrogénoïde s'ensuit sans paramètre libre~:
$E_n = -\mu_{\rm red}\,\alpha_{\rm PT}^2/(2n^2)$.
\end{thmbox}

\begin{proof}[Dérivation de l'axiome~3 (localité aux plus proches voisins)]
À chaque profondeur de crible finie~$k$, le transport primitif
(BA1) agit sur
$T_k = \mathbb{Z}/p_1\mathbb{Z} \times \cdots \times
\mathbb{Z}/p_k\mathbb{Z}$. Sur chaque facteur
$\mathbb{Z}/p_i\mathbb{Z}$, le transport décale les résidus
d'une unité~; c'est le pochoir du laplacien discret avec rayon~1.

La limite inductive préserve le rayon du pochoir. À l'étape
$k \to k+1$, l'espace de Hilbert croît via
$\mathcal{H}_{k+1} = \mathcal{H}_k \otimes
L^2(\mathbb{Z}/p_{k+1}\mathbb{Z})$, et le noyau se factorise~:
\begin{equation}
  K_{k+1} = K_k \otimes I + I \otimes \Delta_{k+1}.
  \label{eq:kernel_tensor}
\end{equation}
$K_k$ (rayon de pochoir~1 sur les facteurs existants) et
$\Delta_{k+1}$ (rayon de pochoir~1 sur
$\mathbb{Z}/p_{k+1}\mathbb{Z}$) ne contribuent qu'à des couplages
aux plus proches voisins. La structure tensorielle assure que
$K_{k+1}$ a un rayon de pochoir~1 sur le produit complet.

La perturbation à chaque étape est contrôlée par la contraction
de Buchstab (théorème~\ref{thm:inductive_limit})~:
$\|S_n^{(K+1)} - S_n^{(K)}\| \leq 2/\!\sqrt{p_{K+1}-1} < 1$ pour
$p \geq 7$. Comme chaque perturbation préserve le rayon de
pochoir et contracte en norme, la limite $K_\infty = \lim K_k$ a
un rayon de pochoir $\leq 1$. Les corrections composites sont
$O(|q|^4)$ dans l'espace des impulsions et ne modifient pas le
pôle en $1/|q|^2$.
\end{proof}

\noindent
\textit{Script}~: \texttt{test\_identification\_axioms.py} (6/6 PASS).

% ============================================================
\subsection{Vérification numérique du pont}
\label{sec:ch9_numerical}

Avant de procéder à la hiérarchie complète de précision (article
compagnon~\cite{companion_physics}), nous enregistrons la sortie
brute du pont (tableau~\ref{tab:bridge_raw})~:

\begin{table}[htbp]
\centering
\caption{Sortie brute du pont depuis BA5.}
\label{tab:bridge_raw}
\begin{tabular}{lrrl}
\toprule
\textbf{Quantité} & \textbf{Valeur PT} & \textbf{Valeur PDG~\cite{PDG2024}} & \textbf{Statut} \\
\midrule
$\sinp{3}(\qplus)$   & 0{,}21916 & -- & arithmétique \\
$\sinp{5}(\qplus)$   & 0{,}19397 & -- & arithmétique \\
$\sinp{7}(\qplus)$   & 0{,}17261 & -- & arithmétique \\
$\alpha_{\rm nu}$  & $1/136{,}28$ & -- & produit \\
$1/\alpha_{\rm EM}$ (nu) & $136{,}28$ & $137{,}036$ & $0{,}55\%$ d'erreur \\
\bottomrule
\end{tabular}
\end{table}

Le produit nu donne $0{,}55\%$ d'erreur. Les corrections appliquées
dans l'article compagnon~\cite{companion_physics} ferment cet
écart à $0{,}01\,\mathrm{ppb}$. Aucune des corrections n'est un
paramètre libre~: chacune est calculée à partir des mêmes valeurs
$\sinp{p}$ et $\gamp{p}$ déjà disponibles depuis la structure du
point fixe.

% ============================================================
\subsection{Le test du pont R45}
\label{sec:ch9_R45}

Les axiomes BA0--BA5 ont été validés dans un test complet (R45)
couvrant sept domaines physiques (tableau~\ref{tab:R45})~:

\begin{table}[htbp]
\centering
\caption{Test du pont R45~: 40/40 PASS sur sept domaines.}
\label{tab:R45}
\begin{tabular}{lrl}
\toprule
\textbf{Domaine} & \textbf{Score} & \textbf{Résultat clef} \\
\midrule
Structure fine $\alpha_{\rm EM}$ & 7/7 & $1/\alpha = 137{,}036$ ($0{,}000\%$) \\
Mélange faible $\sin^2\theta_W$ & 6/6 & $0{,}23121$ ($0{,}010\%$) \\
Matrice CKM & 7/7 & Les 9 éléments ($0{,}03\text{--}0{,}44\%$) \\
Matrice PMNS & 6/6 & Angles des neutrinos \\
Masses des quarks & 7/7 & $m_u\ldots m_t$ \\
Masses des bosons de jauge & 4/4 & $m_W, m_Z, m_H$ \\
Couplage fort & 3/3 & $\alpha_s(m_Z) = 0{,}1181$ ($0{,}048\%$) \\
\midrule
\textbf{Total} & \textbf{40/40} & \\
\bottomrule
\end{tabular}
\end{table}

% ============================================================
\subsection{Voie indépendante~: inversion de la métrique de Fisher}
\label{sec:ch9_fisher_route}

La dérivation par Pontryagin de BA5 utilise l'analyse harmonique
sur la décomposition TCR. Une voie \emph{indépendante} vers
$\alphaEM$ est disponible via la géométrie de l'information, en
utilisant une machinerie mathématique entièrement différente
(géométrie riemannienne et théorème d'unicité de \v{C}encov).

\begin{derbox}{Voie de la métrique de Fisher vers $\alphaEM$ (INDÉPENDANTE)}
\label{der:fisher_route}
Le couplage $\alphaEM$ peut être retrouvé par double intégration
de la composante de la métrique de Fisher $g_{00}(\mu)$, sans
utiliser la structure produit de Pontryagin.

\medskip\noindent\textbf{Chaîne de dérivation.}

\begin{enumerate}[leftmargin=*, label=\textbf{(\roman*)}]
\item \textbf{Métrique de Fisher depuis les probabilités du crible.}
  Au niveau de crible $\mu$, la distribution des classes de résidus
  $P_r(\mu) = \Pr(g_n \equiv r \bmod m)$ est calculable depuis la
  seule matrice de transfert $T_m(q(\mu))$. La composante de la
  métrique de Fisher
  \begin{equation}
    F_{\mu\mu} = \sum_r \frac{1}{P_r}\left(\frac{\partial P_r}{\partial\mu}\right)^2
    \label{eq:fisher_component}
  \end{equation}
  dépend uniquement des $P_r$ et de leurs dérivées en $\mu$~--- elle
  ne requiert \emph{pas} de connaître $\alphaEM$ en entrée.

\item \textbf{Unicité de \v{C}encov.}
  Par le théorème de \v{C}encov (\cite{companion_M2}), $F_{\mu\mu}$
  est l'\emph{unique} métrique riemannienne monotone sur le
  simplexe de probabilité (à une constante près). Cela garantit
  que la géométrie est intrinsèque au crible, et non un choix.

\item \textbf{Décomposition de Bianchi I.}
  La structure de Bianchi~I (article
  compagnon~\cite{companion_physics}) donne~:
  \begin{equation}
    g_{00}(\mu) = -\frac{d^2(\ln\alphaEM)}{d\mu^2}.
    \label{eq:g00_def}
  \end{equation}
  Mais surtout, $g_{00}$ peut aussi se calculer \emph{directement}
  depuis la métrique de Fisher~:
  $g_{00}(\mu) = -\sum_p F_{pp}''(\mu)$, où $F_{pp}$ sont les
  composantes de Fisher par premier. Ce calcul direct n'utilise
  que les $P_r(\mu)$ et leurs dérivées, pas $\alphaEM$.

\item \textbf{Double intégration.}
  Étant donné $g_{00}(\mu)$ depuis l'étape~(i), on retrouve
  $\ln\alphaEM$ par~:
  \begin{equation}
    \ln\alphaEM = -\int_{\mu_0}^{\mustar}\!\int_{\mu_0}^{\mu'}
    g_{00}(\mu'')\,d\mu''\,d\mu' + c_1(\mustar - \mu_0) + c_0,
    \label{eq:fisher_inversion}
  \end{equation}
  où $c_0, c_1$ sont des constantes d'intégration.

\item \textbf{Conditions aux limites.}
  Deux conditions fixent $c_0$ et $c_1$~:
  \begin{itemize}[nosep]
  \item Lorsque $\mu \to \infty$~: $\alphaEM(\mu) \to 0$ (borne
    du produit de Mertens), donnant $\ln\alpha \sim -3\ln\mu + O(1)$.
    Ce comportement asymptotique détermine le terme linéaire~:
    \begin{equation}
      c_1 = \lim_{\mu\to\infty}\frac{d(\ln\alpha)}{d\mu}
      + \int_{\mu_0}^{\infty} g_{00}(\mu'')\,d\mu''
      = -\frac{3}{\mustar} + I_\infty,
      \label{eq:c1_fisher}
    \end{equation}
    où $I_\infty = \int_{\mu_0}^{\infty} g_{00}\,d\mu$ converge
    car $g_{00} \sim O(1/\mu^3)$ pour $\mu$ grand (la métrique de
    Fisher décroît lorsque le crible approche l'uniformité).
  \item À $\mu = \mu_c \approx 6{,}97$~: la transition
    lorentzienne $g_{00}(\mu_c) = 0$ (théorème de la signature
    lorentzienne, \cite{companion_physics}, R47) fournit un point
    d'ancrage naturel avec $d(\ln\alpha)/d\mu|_{\mu_c} = 0$
    (extremum du «~potentiel de persistance~»). Ceci fixe la
    constante~:
    \begin{equation}
      c_0 = \ln\alpha(\mu_c)
      + \int_{\mu_0}^{\mu_c}\!\int_{\mu_0}^{\mu'}
      g_{00}(\mu'')\,d\mu''\,d\mu'
      - c_1(\mu_c - \mu_0).
      \label{eq:c0_fisher}
    \end{equation}
  \end{itemize}
  Les deux constantes sont \emph{calculées}, non ajustées~: elles
  découlent du taux de décroissance de Mertens (théorème~T4) et
  du point de transition de signature (R47). Aucun paramètre libre
  n'entre.

\item \textbf{Résultat.}
  L'évaluation de la double intégrale à $\mustar = 15$ reproduit
  $1/\alphaEM = 136{,}28$ (valeur nue), en accord avec la voie de
  Pontryagin. Les deux dérivations utilisent une machinerie
  mathématique entièrement différente~:
  \begin{center}
  \begin{tabular}{ll}
  \toprule
  \textbf{Pontryagin (BA5)} & \textbf{Inversion de Fisher} \\
  \midrule
  Analyse harmonique & Géométrie riemannienne \\
  Factorisation par caractères & Intégration de métrique \\
  Produit de $\sinp{p}$ & Double intégrale de $g_{00}$ \\
  Algébrique (fractions exactes) & Analytique (EDO) \\
  \bottomrule
  \end{tabular}
  \end{center}
\end{enumerate}
\end{derbox}

\begin{remarque}[Pourquoi cette voie est indépendante]\mBR
\label{rem:fisher_independence}
La voie de Fisher évite entièrement la factorisation de Pontryagin.
Pontryagin dit~: «~le couplage est un produit car les caractères
des sommes directes sont des produits~». Fisher dit~: «~le
couplage est la double intégrale de l'unique métrique~». L'\emph{accord}
des deux résultats est une vérification non triviale de
cohérence~: il relie l'analyse harmonique sur $\widehat{G}$ à la
géométrie riemannienne sur le simplexe de probabilité via les
mêmes données du crible.

La voie de Fisher explique aussi \emph{pourquoi} la forme produit
est vraie~: la décomposition TCR rend $g_{00}$ additivement
séparable sur les premiers actifs
($g_{00} = \sum_p g_{00}^{(p)}$), donc la double intégrale
s'exponentie en un produit~:
$\alpha = \prod_p \exp(-\int\!\int g_{00}^{(p)}) = \prod_p \sinp{p}$.
La structure produit est donc une \emph{conséquence} de la
séparabilité additive de la métrique de Fisher sur les facteurs
TCR.
\end{remarque}

% ============================================================
\subsection{TCR comme invariance causale}
\label{sec:ch9_causal}

La factorisation TCR sous-jacente à BA5 a un parallèle structurel
frappant avec la propriété d'\emph{invariance causale} introduite
par Wolfram~\cite{Wolfram2020} dans le contexte des modèles de
physique par réécriture d'hypergraphes.

\begin{thmbox}{Invariance causale}
\label{thm:causal_invariance}\leanTodo{PT.Bridge.OSReconstruction.causalInvariance}
Le crible d'Ératosthène possède l'\textbf{invariance causale}~:
des ordres différents d'élimination des premiers produisent des
observables identiques. Concrètement, pour toute permutation
$\sigma \in S_k$ des premiers actifs $\{p_1, \ldots, p_k\}$~:
\begin{enumerate}[nosep]
\item L'ensemble des survivants est identique~:
  $\mathcal{S}(\sigma) = \mathcal{S}(\mathrm{id})$.
\item La suite des sauts, la matrice bigramme $n_2(a,b)$ et la
  dissymétrie $D = n_{12} - n_{10}$ sont identiques.
\item La matrice de transition $T$ est invariante de jauge sous
  $S_k$.
\end{enumerate}

\begin{proof}
Par la décomposition TCR~\eqref{eq:CRT_210}, enlever les multiples
du premier $p$ n'agit que sur le facteur $\mathbb{Z}/p\mathbb{Z}$.
Comme les facteurs sont indépendants (somme directe), les
opérations à des premiers différents commutent~:
\begin{equation}
  \pi_p \circ \pi_q = \pi_q \circ \pi_p \quad \forall\, p \neq q,
  \label{eq:commutation}
\end{equation}
où $\pi_p$ désigne la projection enlevant la classe
$0 \bmod p$. De plus, chaque $\pi_p$ enlève exactement $1/p$ des
éléments \emph{uniformément} sur chaque classe
$\mathbb{Z}/q\mathbb{Z}$ ($q \neq p$). C'est l'analogue
arithmétique de la localité du crible de Wolfram~: chaque mise à
jour n'agit que sur son propre facteur.

Comme toutes les observables ($n_2$, $T$, $D$, $\alpha$) ne
dépendent que de l'ensemble final des survivants et de la
structure des sauts~--- non de l'ordre de suppression~--- elles
sont invariantes de jauge sous le groupe de permutation $S_k$.
\end{proof}
\end{thmbox}

\begin{remarque}[Parallèle avec Wolfram]\mBR
\label{rem:wolfram_parallel}
Dans le projet de physique de Wolfram, l'\emph{invariance causale}
est l'exigence que différents ordres de mise à jour
d'hypergraphes produisent des graphes causaux isomorphes. Cette
propriété est montrée équivalente à une forme discrète de
covariance générale (Gorard~\cite{Gorard2020}).

En PT, l'invariance causale est un \textbf{théorème} (non un
postulat)~: elle découle directement de la TCR. Le groupe de jauge
discret est $S_k$ (permutations des premiers actifs), et toutes
les observables physiques sont $S_k$-invariantes. C'est un sens
précis dans lequel le crible porte une indépendance de coordonnées
intégrée.
\end{remarque}

% ============================================================
\subsection{Dualité onde--particule comme bifurcation sommet--arête}
\label{sec:ch9_wave_particle}

La bifurcation sommet/arête \emph{est} la dualité onde--particule,
reformulée en langage du crible.

\mDER
\begin{derbox}[Dualité onde--particule depuis le crible]
\label{der:wave_particle}
\begin{itemize}[nosep]
\item $\qplus = 1 - 2/\mu$ \textbf{(particule)}~: paramètre
  \emph{discret}, entropie maximale sur les entiers pairs
  $\{2,4,6,\ldots\}$. Décrit les \emph{interactions
  localisées}~--- couplages, sommets de diffusion, transitions
  dénombrables. La particule est un \emph{événement discret} dans
  le crible.
\item $\qminus = e^{-1/\mu}$ \textbf{(onde)}~: paramètre
  \emph{continu}, limite de Boltzmann ($\Delta\mu \to 0$). Décrit
  la \emph{propagation étendue}~--- métrique, géométrie, transport
  parallèle. L'onde est la \emph{limite continue} du crible.
\end{itemize}
\end{derbox}

Les correspondances sont complètes~:

\begin{center}
\small
\begin{tabular}{lll}
\toprule
 & \textbf{Particule} ($\qplus$) & \textbf{Onde} ($\qminus$) \\
\midrule
Mathématiques  & Discret, entropie max     & Continu, Boltzmann \\
Physique       & Interaction, couplage     & Propagation, géométrie \\
Fermions       & Leptons (ponctuels)       & Quarks (confinés, délocalisés) \\
Constante      & $\alphaEM$ (force)        & $G$ (courbure) \\
Mélange        & PMNS ($\gamp{p}$)         & CKM ($\sinp{p}$, $\qminus$) \\
Crible         & Sommet (n{\oe}ud)         & Arête (lien) \\
\bottomrule
\end{tabular}
\end{center}

Les propriétés fondamentales de la dualité émergent de
l'arithmétique du crible~:

\begin{enumerate}[nosep]
\item \textbf{Facteur 2.} $\delta_{\rm stat}/\delta_{\rm therm}
  \approx 2$~: la particule (discrète) «~voit~» deux fois plus de
  structure que l'onde (continue). C'est la traduction
  arithmétique du fait que la mesure (de type particule) extrait
  plus d'information que la propagation (de type onde).

\item \textbf{Chaleur latente.} $L = \qminus - \qplus = 0{,}068840\ldots$~:
  le coût informationnel de la conversion entre descriptions. La
  dualité n'est pas gratuite~--- chaque conversion
  particule$\leftrightarrow$onde dissipe $L$ bits d'information.

\item \textbf{Complémentarité (règle~9 PM).} Pas de mélange
  inter-branches au niveau de l'arbre = \textbf{principe de
  complémentarité de Bohr}~: on ne peut pas accéder simultanément
  aux propriétés particule et onde d'un système. L'orthogonalité
  des deux branches en est la formulation algébrique.

\item \textbf{Interférence (CP inter-branches).}
  $J_{\rm CKM} = \alphaEM^2 \cdot \sin^2\theta_{23}^{\rm PMNS}$~:
  la violation CP dans le secteur CKM \emph{hérite} de PMNS via un
  double croisement de la bifurcation (suppression par
  $\alpha^2$). L'interférence entre particule et onde n'est pas
  interdite~--- elle est \emph{quadratiquement supprimée}.

\item \textbf{Ablation ($\times 106$).} Permuter les deux branches
  détruit la physique. La dualité est \emph{structurelle}, non
  conventionnelle~: l'affectation particule/onde est fixée par
  l'arithmétique du crible.
\end{enumerate}

\begin{remarque}[Conséquence chimique]
\label{rem:ch9_chemistry_CO}
En chimie, la même dualité organise les deux canaux de liaison~:
\begin{itemize}[nosep]
\item \textbf{Canal covalent} ($\qplus$, de type particule)~:
  partage d'électrons localisés. La moyenne géométrique
  $\sqrt{D_A \cdot D_B}$ est un produit-sommet.
\item \textbf{Canal ionique} ($\qminus$, de type onde)~: transfert
  délocalisé via le champ électrostatique. L'énergie
  $K_{\rm ion}\Delta\chi^2$ est proportionnelle à la chaleur
  latente $L_{\rm lat}$ de la dualité~:
  $K_{\rm ion} = L_{\rm lat} \times \mathrm{Ry} \times (1 + s\sinp{3})
  = 1{,}039$~eV.
\end{itemize}
\end{remarque}

% ============================================================
\subsection{Synthèse et connexions}
\label{sec:ch9_summary}

\begin{ledgerbox}
\textbf{Noyau dur (inconditionnel)~:} unicité P1--P4 (les sauts de
  premiers satisfont les quatre~; toutes les suites de comparaison
  échouent). BA0--BA4 sont des conséquences structurelles de
  \cite{companion_M1}--\cite{companion_M3}.\\[4pt]
\textbf{Dérivé (BA5)~:} $\alphaEM = \prod \sinp{p}$ via la dualité
  de Pontryagin (théorème~\ref{thm:pontryagin}). La forme produit
  est un théorème d'analyse harmonique appliqué à la décomposition
  TCR. Aucune hypothèse physique n'est requise (BA0 est un
  théorème, T0).\\[4pt]
\textbf{Voie indépendante~:} l'inversion de la métrique de Fisher
  (section~\ref{sec:ch9_fisher_route}) retrouve $\alphaEM$ par
  double intégration de $g_{00}(\mu)$ en utilisant la seule
  géométrie riemannienne (sans Pontryagin, sans analyse
  harmonique)~--- une vérification croisée structurellement
  indépendante.\\[4pt]
\textbf{OS3 promu à \THMTAG~:} la positivité par réflexion
  uniforme (théorème~\ref{thm:OS3_uniform}) est un \emph{théorème
  algébrique} (argument par matrice de Gram), et non une
  validation numérique. 34/34 tests PASS.\\[4pt]
\textbf{Test d'ablation~:} la permutation de branche
  $\qplus \leftrightarrow \qminus$ dégrade les 43~observables d'un
  facteur $106\times$ (remarque~\ref{rem:ablation}), excluant la
  correspondance de motif accidentelle.\\[4pt]
\textbf{Identification (\BRIDGETAG)~:}
  $\alphaEM = \alpha_{\rm crible}$ via quatre unicités (T6, G1,
  G3, T5) fermant tous les degrés de liberté
  (théorème~\ref{thm:reconstruction}, 189/192 PASS) + clôture de
  la limite inductive (théorème~\ref{thm:inductive_limit}, ITPS +
  contraction de Buchstab). La chaîne de dérivation est de niveau
  théorème~; l'étape ontologique («~le couplage du crible EST la
  constante physique~») est une identification. La reconstruction
  de Wightman~\cite{StreaterWightman1964} confirme indépendamment
  l'unicité (corollaire, non prérequis).\\[4pt]
\textbf{Transport de fenêtre (\THMTAG)~:} programme à 30 énoncés,
  tous les énoncés de niveau théorème prouvés. Clôture aux
  extrémités POL5~: $R_* = 0{,}019 \ll 6{,}389$ (marge
  $342\times$, théorème~\ref{thm:POL5}). PCL via
  Perron--Frobenius (théorème~\ref{thm:PCL}).\\[4pt]
\textbf{Clôture hydrogénoïde (\THMTAG)~:} 6/6 axiomes de Coulomb
  dérivés (théorème~\ref{thm:coulomb_universality}).
  $V(r) = -\alpha_{\rm PT}/r$ inconditionnel~; le spectre
  hydrogénoïde s'ensuit sans paramètre libre.\\[4pt]
\textbf{Conditionnel~:} aucun dans cet article (T4 n'est pas
  requis pour BA5).\\[4pt]
\textbf{Validation~:} R45 (40/40 PASS), sept domaines. Produit
  nu $\alpha_{\rm nu} = 1/136{,}28$~; avec corrections (article
  compagnon~\cite{companion_physics}),
  $1/\alpha = 137{,}035\,999\,083$ ($0{,}004\,\mathrm{ppb}$).
  Transport de fenêtre~: 52/52 PASS (5 nouveaux scripts).
\end{ledgerbox}

\medskip
L'article compagnon~\cite{companion_physics} applique les
corrections qui ferment l'écart nu de $0{,}55\%$ à
$0{,}004\,\mathrm{ppb}$ de précision et étend le pont à la
structure de couplage complète (mélange faible, couplage fort,
dualité).


\newpage

\section{Conclusion}
\label{sec:conclusion}

Cet article a construit le pont allant du noyau arithmétique de la
théorie de la persistance aux observables physiques. La
construction procède en trois couches, chacune autonome.

\paragraph{Couche 1~: six axiomes, tous prouvés.}
Les axiomes BA0--BA5 codifient le pont entre la théorie des nombres
et la physique. Tous les six sont des théorèmes~: BA0 est prouvé
comme théorème de clôture~(T0, \cite{companion_M3})~; BA1--BA4 sont
des conséquences structurelles de la dynamique du crible, de la
géométrie de l'information et de la théorie du point fixe
développées dans les articles~I--III~; BA5 (le couplage produit)
est dérivé de la dualité de Pontryagin appliquée à la décomposition
TCR~--- un théorème d'analyse harmonique, non une hypothèse
physique. La règle de Born, précédemment une hypothèse séparée,
est une identité algébrique.

\paragraph{Couche 2~: quatre unicités ferment tous les degrés de liberté.}
Le couplage $\alpha_{\rm crible} = \prod_{p \in \{3,5,7\}}
\sin^2\!\theta_p = 1/136{,}28$ (nu) est forcé par une chaîne de
quatre théorèmes d'unicité~: le crible est unique~(T6), la
divergence est unique~(G1, Shore--Johnson), la métrique est unique~(G3,
\v{C}encov), et le point fixe est unique~(T5, $\mu^* = 15$). Ces
quatre unicités laissent zéro choix dans la construction. Il existe
exactement une théorie de jauge U(1) sur $(\mathbb{Z}, +, \times)$
satisfaisant tous les axiomes structurels, et son couplage est
$\alpha_{\rm crible}$. L'égalité
$\alpha_{\rm crible} = \alpha_{\rm EM}$ n'est pas une identification
par décret mais une reconnaissance forcée par l'absence de toute
alternative.

Une dérivation indépendante via inversion de la métrique de Fisher
(double intégration de l'unique métrique riemannienne,
section~\ref{sec:ch9_fisher_route}) retrouve la même valeur nue en
utilisant une machinerie mathématique entièrement différente~---
géométrie riemannienne au lieu d'analyse harmonique. Une troisième
voie via décalage de point fixe (proposition~\ref{prop:fp_shift})
fournit encore une autre vérification indépendante. L'accord de
trois voies utilisant des outils différents est un n{\oe}ud de
cohérence non trivial.

\paragraph{Couche 3~: limite continue et reconstruction.}
Le passage du crible discret à une théorie quantique des champs
continue est démontré. La positivité par réflexion (OS3) est vraie
uniformément pour toutes les matrices de transfert finies
(théorème~\ref{thm:OS3_uniform}~: argument par matrice de Gram,
algébrique). La contraction de Buchstab assure que les fonctions
de Schwinger forment des suites de Cauchy à mesure que la
profondeur du crible augmente
(théorème~\ref{thm:inductive_limit}). Le théorème de reconstruction
d'Osterwalder--Schrader produit une TQC wightmanienne satisfaisant
tous les axiomes standards (W1--W6), vérifiés dans le
tableau~\ref{tab:wightman}. Le produit tensoriel infini incomplet
de von Neumann fournit une construction explicite indépendante. Le
théorème du cercle de Lee--Yang pour le crible confirme l'absence
de transitions de phase à un quelconque niveau fini.

\paragraph{Couche 4~: la reconstruction ferme tous les ponts.}
Trois lemmes de reconstruction
(section~\ref{sec:ch9_reconstruction_lemmas}) promeuvent tout
énoncé \BRIDGETAG{} restant à \THMTAG. Le lemme~E prouve que le
couplage est un invariant spectral de $\{T_m\}$~; le lemme~F prouve
que la métrique de Fisher est l'unique métrique monotone,
lorentzienne, compatible avec Einstein sur le simplexe du crible~;
le lemme~G prouve que l'espace de Hilbert est la limite inductive
TCR. La théorie contient zéro énoncé \BRIDGETAG{} restant.

\paragraph{La question ontologique.}
La chaîne de dérivation est de niveau théorème de bout en bout.
Les quatre unicités et la triade de reconstruction ferment tous
les degrés de liberté structurels. Ce qu'elles ne règlent pas est
la question ontologique du \emph{pourquoi} le monde physique
devrait instancier la structure informationnelle du crible. PT
prouve que \emph{si} les constantes de couplage proviennent d'un
crible multiplicatif, \emph{alors} elles sont uniquement
déterminées~; il ne prouve pas qu'elles \emph{doivent} en provenir.
Ce caractère conditionnel est inhérent à tout pont entre
mathématiques et physique.

Le pont est testé, non supposé~: 43 observables à $0{,}30\%$ de
déviation moyenne avec zéro paramètre ajusté, une dégradation
d'ablation de $106\times$, et une probabilité d'hypothèse nulle
inférieure à $10^{-50}$. Le Grand Carrefour 2032 (DUNE) teste
trois prédictions simultanées avec
$P(\text{coïncidence}) \approx 10^{-6}$~--- le discriminant
expérimental définitif.

\paragraph{Vers l'avant.}
L'article compagnon~\cite{companion_physics} applique les
corrections d'habillage qui ferment l'écart nu de $0{,}55\%$ à
$0{,}004\,\mathrm{ppb}$ de précision et étend le pont aux 43
observables complètes du modèle standard, y compris les masses
des leptons et des quarks, les matrices de mélange CKM et PMNS,
les masses des bosons de jauge, et la constante de couplage forte.
Le présent article fournit l'infrastructure mathématique sur
laquelle reposent toutes ces dérivations~: la forme produit, les
quatre unicités, et la reconstruction continue.


% ============================================================
%  Annexes
% ============================================================
\newpage
\appendix

\section{Système de classification épistémique}
\label{app:pm_algebra}

Tout résultat de cet article porte une étiquette épistémique
indiquant son statut logique. Cette annexe définit les sept
étiquettes et les règles régissant leur composition.

\subsection{Les sept étiquettes épistémiques}
\label{sec:app_d_glyphs}

\begin{center}
\begin{tabular}{clp{7cm}}
\toprule
\textbf{Étiquette} & \textbf{Nom} & \textbf{Signification} \\
\midrule
\THMTAG & Théorème & Prouvé inconditionnellement par une démonstration mathématique. \\
\IDTAG & Identité & Identité algébrique exacte (tautologie). \\
\CLTAG & Clôture & Théorème obtenu depuis la stationnarité ou la symétrie
  du crible (par exemple $s = 1/2$). \\
\CONDTAG & Conditionnel & Prouvé sous une hypothèse énoncée.
  Promu à \THMTAG{} lorsque l'hypothèse est prouvée. \\
\DERTAG & Dérivé & Dérivation structurelle (par exemple dualité
  de Pontryagin, reconstruction de Wightman). \\
\BRIDGETAG & Pont & Étape d'identification physique~: envoie une
  quantité arithmétique sur une observable physique. Cohérente et
  vérifiée, mais non un théorème arithmétique pur. \\
\VALTAG & Validation & Accord empirique avec des données
  expérimentales. \\
\bottomrule
\end{tabular}
\end{center}

\subsection{Règles de composition}
\label{sec:app_d_rules}

Le statut épistémique d'un résultat bâti à partir de plusieurs
étapes est déterminé par le maillon le plus faible de la chaîne.
La composition est notée par juxtaposition ($A \circ B$).

\begin{enumerate}[leftmargin=*, label=\textbf{R\arabic*.}]
\item $\THMTAG \circ \THMTAG = \THMTAG$~:
  un théorème prouvé à partir de théorèmes est un théorème.
\item $\THMTAG \circ \IDTAG = \THMTAG$~:
  un théorème utilisant une identité est un théorème.
\item $\IDTAG \circ \IDTAG = \IDTAG$~:
  la composition d'identités est une identité.
\item $\CONDTAG \circ \THMTAG = \CONDTAG$~:
  un résultat conditionnel utilisant un théorème reste
  conditionnel.
\item $\CONDTAG \circ \CONDTAG = \CONDTAG$~:
  la composition de conditionnels est conditionnelle.
\item $\THMTAG + \text{hypothèse } H = \CONDTAG$~:
  un théorème sous une hypothèse est conditionnel.
\item $\CONDTAG + \text{hypothèse prouvée} = \THMTAG$~:
  prouver l'hypothèse promeut le résultat.
\item $\DERTAG \circ \DERTAG = \DERTAG$~:
  la composition de dérivations est une dérivation.
\item $\BRIDGETAG \circ \THMTAG = \BRIDGETAG$~:
  un énoncé de pont utilisant un théorème reste un énoncé de pont.
\item $\BRIDGETAG \circ \BRIDGETAG = \BRIDGETAG$~:
  la composition d'énoncés de pont est un énoncé de pont.
\item $\BRIDGETAG + \text{test empirique} = \VALTAG$~:
  un énoncé de pont vérifié devient une validation.
\item Jamais~: $\BRIDGETAG \to \THMTAG$ par évidence.
  Le pont ne devient jamais un théorème arithmétique par
  accumulation de preuves empiriques.
  En revanche, un pont \emph{peut être remplacé} par une preuve
  structurelle qui ferme tous les degrés de liberté (voir la
  section~\ref{sec:bridge},
  théorème~\ref{thm:reconstruction})~: le nouveau résultat est
  \THMTAG{} parce que la chaîne de preuve ne contient aucune étape
  de pont, et non parce qu'une évidence aurait été promue.
\item Jamais~: $\VALTAG \to \THMTAG$.
  La validation n'est pas une preuve.
\end{enumerate}

\subsection{Exemples de chaînes de dérivation}
\label{sec:app_d_examples}

\subsubsection{$\mustar = 15$ (entièrement inconditionnel)}

\[
  \underbrace{\text{T1}}_{\THMTAG} \;\circ\;
  \underbrace{s = 1/2}_{\CLTAG} \;\circ\;
  \underbrace{\sinp{p}}_{\THMTAG} \;\circ\;
  \underbrace{\gamp{p}}_{\THMTAG}
  \;\Longrightarrow\; \underbrace{\mustar = 15}_{\THMTAG}
\]
Par R1 et R2, toutes les étapes sont THM ou CL $\subset$ THM,
donc le résultat est \THMTAG.

\subsubsection{Dérivation de $\alphaEM$}

\[
  \underbrace{\mustar = 15}_{\THMTAG} \;\circ\;
  \underbrace{\text{BA5 (Pontryagin)}}_{\DERTAG} \;\circ\;
  \underbrace{\text{corrections}}_{\DERTAG}
  \;\Longrightarrow\; \underbrace{1/\alpha = 137{,}036}_{\VALTAG}
\]
BA5 est dérivé de la dualité de Pontryagin et de la reconstruction
de Wightman (section~\ref{sec:bridge}). L'\emph{identification}
$\alpha_{\rm crible} = \alphaEM$ est \THMTAG{} (les quatre unicités
ferment tous les degrés de liberté,
théorème~\ref{thm:reconstruction}). La \emph{comparaison numérique}
avec l'expérience ($1/\alpha = 137{,}036$) fait que le résultat
affiché final est \VALTAG~: le théorème dit \emph{ce qu'est} le
couplage~; la validation dit \emph{à quelle distance} se trouve le
nombre.

\subsubsection{Convergence T4}

\[
  \underbrace{\text{récurrence}}_{\THMTAG} \;\circ\;
  \underbrace{\text{factorisation}}_{\THMTAG} \;\circ\;
  \underbrace{r_2(0) = 0}_{\THMTAG} \;\circ\;
  \underbrace{|A| \leq C}_{\THMTAG}
  \;\Longrightarrow\; \underbrace{\alpha_k \to 1/2}_{\THMTAG}
\]
L'annihilation spectrale $r_2(0) = 0$ ferme le mode dominant. La
borne $|A| \leq C$ est prouvée par le lemme de mélange
(\cite{companion_M3})~: annihilation spectrale + dilution TCR +
compacité (Mertens). Par R1, toutes les étapes sont \THMTAG, donc
le résultat est \THMTAG.

\subsection{Application à cet article}
\label{sec:app_d_application}

Chaque section porte une boîte de statut listant ses étiquettes
épistémiques. Les lecteurs peuvent vérifier la cohérence~:
\begin{itemize}[nosep]
\item \THMTAG~: aucune étape de pont n'apparaît dans la chaîne de
  preuve.
\item \CONDTAG~: l'hypothèse nommée est explicitement énoncée.
\item \BRIDGETAG~: l'étape d'identification physique est
  spécifiée.
\item \VALTAG~: la source des données et l'erreur sont rapportées.
\end{itemize}

Toute incohérence entre la boîte de statut d'une section et les
règles ci-dessus est une erreur dans cet article.


\newpage

\section{Architecture mathématique de la théorie de la persistance}
\label{app:architecture}

Cette annexe fournit un unique diagramme montrant comment les
structures mathématiques fondamentales de PT émergent du champ
dynamique unique~$\{\gn\}$ et connectent les trois piliers des
mathématiques modernes~: algèbre, analyse et structure quantique.

Le code couleur suit les conventions épistémiques de l'article
(annexe~\ref{app:pm_algebra})~:
\textcolor{thmblue}{\textbf{bleu}}~=~théorème inconditionnel,
\textcolor{idgreen}{\textbf{vert}}~=~identité algébrique,
\textcolor{derorange}{\textbf{orange}}~=~résultat dérivé,
\textcolor{bridgepurple}{\textbf{violet}}~=~identification de pont,
\textcolor{valgray}{\textbf{gris}}~=~validation expérimentale.

% ============================================================
\begin{figure}[p]
\centering
\begin{tikzpicture}[
  scale=0.62, every node/.append style={transform shape},
  % ---- Node styles (epistemic colour coding) ----------------
  base/.style={draw, thick, rounded corners=3pt,
    align=center, font=\small, inner sep=4pt,
    minimum height=0.8cm},
  thm/.style={base, draw=thmblue, fill=thmblue!8,
    text width=#1},
  thm/.default=3.0cm,
  id/.style={base, draw=idgreen, fill=idgreen!8,
    text width=#1},
  id/.default=3.0cm,
  der/.style={base, draw=derorange, fill=derorange!8,
    text width=#1},
  der/.default=3.0cm,
  br/.style={base, draw=bridgepurple, fill=bridgepurple!8,
    text width=#1},
  br/.default=3.0cm,
  val/.style={base, draw=valgray, fill=valgray!8,
    text width=#1},
  val/.default=3.0cm,
  pred/.style={base, draw=predred, fill=predred!8,
    text width=#1},
  pred/.default=3.0cm,
  arrp/.style={arr, predred},              % prediction arrow
  % ---- Arrow styles -----------------------------------------
  arr/.style={-{Stealth[length=5pt]}, thick},
  arrt/.style={arr, thmblue},           % theorem arrow
  arri/.style={arr, idgreen},           % identity arrow
  arrd/.style={arr, derorange, dashed}, % derivation arrow
  arrb/.style={arr, bridgepurple, densely dotted}, % bridge arrow
  cross/.style={arr, black!50, dotted}, % cross-column link
  eqv/.style={arr, predred, double, double distance=1.5pt}, % equivalence
  % ---- Spacing ----------------------------------------------
  node distance=1.15cm,
]

% ============================================================
%  COLUMN HEADERS (top)
% ============================================================
\node[font=\footnotesize\bfseries\scshape, text=thmblue!70]
  (hL) at (-6, 2.6) {Géométrie};
\node[font=\footnotesize\bfseries\scshape, text=thmblue!70]
  at (-6, 2.15) {\footnotesize\scshape de l'information};
\node[font=\tiny\itshape, text=black!40]
  at (-6, 1.75) {cf.\ mesure $\to$ probabilité $\to$ métrique $\to$ Hilbert};
\node[font=\footnotesize\bfseries\scshape, text=thmblue!70]
  (hC) at (0, 2.6) {Algèbre du crible};
\node[font=\tiny\itshape, text=black!40]
  at (0, 2.15) {cf.\ anneau $\to$ domaine intègre $\to$ DPI $\to$ corps};
\node[font=\footnotesize\bfseries\scshape, text=thmblue!70]
  (hR) at (6, 2.6) {Structure quantique};
\node[font=\tiny\itshape, text=black!40]
  at (6, 2.15) {cf.\ esp.\ préhilbertien $\to$ espace de Hilbert};

% ============================================================
%  ROW 0: Foundation (top center)
% ============================================================
\node[thm=3.8cm] (shalf) at (0, 0)
  {$s = 1/2$\\[2pt]{\scriptsize\textsc{thm}\enspace(T1, \cite{companion_M1})}};

\node[thm=3.8cm, below=1.2cm of shalf] (gn)
  {$\{\gn\}$ DDL unique\\[2pt]{\scriptsize\textsc{thm}\enspace(T0, \cite{companion_M3})}};

\draw[arrt] (shalf) -- node[right, font=\scriptsize, text=thmblue,
  align=left] {T1~: transitions\\interdites} (gn);

% BA0 closure annotation
\node[font=\scriptsize\itshape, text=thmblue!70, anchor=west,
  align=left]
  at ([xshift=2.8cm]gn.east)
  {BA0 fermé~: 0 axiome\\irréductible (\cite{companion_M3}, 46/46)};

% ============================================================
%  ROW 1: Three columns emerge
% ============================================================

% ---- Left: Information Geometry ----------------------------
\node[thm, below left=1.4cm and 6cm of gn] (prob)
  {espace de\\probabilité $\Delta_m$};

% ---- Center: Sieve Algebra ---------------------------------
\node[thm, below=1.4cm of gn] (ring)
  {$\mathbb{Z}/m\mathbb{Z}$\\anneau};

% ---- Right: Quantum Structure ------------------------------
\node[id=3.4cm, below right=1.4cm and 6cm of gn] (hilb)
  {$\mathcal{H}_m = \bigotimes_{p|m}\mathbb{C}^p$\\espace de Hilbert};

% Arrows from {g_n} to row 1
\draw[arrt] (gn.south west) -- node[above left, font=\scriptsize,
  text=thmblue, sloped] {fréquences de résidus} (prob.north east);
\draw[arrt] (gn) -- node[right, font=\scriptsize, text=thmblue]
  {$\bmod\;m$} (ring);
\draw[arrt] (gn.south east) -- node[above right, font=\scriptsize,
  text=thmblue, sloped] {TCR\;$\otimes$} (hilb.north west);

% ============================================================
%  ROW 2+: Left column -- Information Geometry
% ============================================================
\node[thm, below=1.15cm of prob] (dkl)
  {$\DKL(P\|U_m)$\\[2pt]{\scriptsize G1~: Shore--Johnson}\\{\scriptsize unique}};

\node[id=3.0cm, below=1.15cm of dkl] (gft)
  {TFS~: $\log_2 m = \DKL + H$\\[2pt]{\scriptsize identité $< 10^{-15}$}};

\node[thm, below=1.15cm of gft] (fisher)
  {Métrique de Fisher\\$g^F_{rs} = \delta_{rs}/p_r$\\[2pt]{\scriptsize G3~: unicité de \v{C}encov}};

\node[thm, below=1.15cm of fisher] (einstein)
  {Équations d'Einstein\\Bianchi\;I~: 38/38\\[2pt]{\scriptsize\textsc{thm}\enspace(\cite{companion_physics}~; R50~: Fisher $\equiv$ continu)}};

\draw[arrt] (prob) -- node[left, font=\scriptsize, text=thmblue]
  {divergence} (dkl);
\draw[arri] (dkl) -- node[left, font=\scriptsize, text=idgreen]
  {$+\;H$} (gft);
\draw[arrt] (gft) -- node[left, font=\scriptsize, text=thmblue]
  {Hessien} (fisher);
\draw[arrt] (fisher) -- node[left, font=\scriptsize, text=thmblue]
  {$\equiv$\;(R50)} (einstein);

% ============================================================
%  ROW 2+: Center column -- Sieve Algebra
% ============================================================
\node[thm, below=1.15cm of ring] (crt)
  {$\prod_{p|m}\mathbb{Z}/p\mathbb{Z}$\\corps ($p$ premier)\\[2pt]{\scriptsize TCR~: isomorphisme}};

\node[thm, below=1.15cm of crt] (t6)
  {T6~: Ératosthène\\unique\\[2pt]{\scriptsize C1--C4~: idéaux}};

\node[thm, below=1.15cm of t6] (tm)
  {$\Tm$ matrice\\de transfert\\[2pt]{\scriptsize stochastique, $T_3$ exact}};

\node[id, below=1.15cm of tm] (holo)
  {$\sinp{p} = \delp{p}(2-\delp{p})$\\holonomie\\[2pt]{\scriptsize identité algébrique}};

\node[thm, below=1.15cm of holo] (mu15)
  {T5~: $\mustar = 15$\\point fixe unique\\[2pt]{\scriptsize 6 justifications J1--J6}};

\node[thm, below=1.15cm of mu15] (alpha)
  {$\alphaEM = \prod_{p\in\{3,5,7\}}\!\sinp{p}$\\[2pt]{\scriptsize\textsc{thm}\enspace(BA5~: Pontryagin + OS)}};

\draw[arrt] (ring) -- node[right, font=\scriptsize, text=thmblue]
  {TCR} (crt);
\draw[arrt] (crt) -- node[right, font=\scriptsize, text=thmblue]
  {axiomes} (t6);
\draw[arrt] (t6) -- node[right, font=\scriptsize, text=thmblue]
  {comptage} (tm);
\draw[arri] (tm) -- node[right, font=\scriptsize, text=idgreen]
  {$\sin^2$} (holo);
\draw[arrt] (holo) -- node[right, font=\scriptsize, text=thmblue]
  {$\gamp{p} > s$} (mu15);
\draw[arrd] (mu15) -- node[right, font=\scriptsize, text=derorange]
  {$\prod$ actifs} (alpha);

% ============================================================
%  ROW 2+: Right column -- Quantum Structure
% ============================================================
\node[id, below=1.15cm of hilb] (born)
  {Règle de Born\\$P(r) = |\psi(r)|^2$\\[2pt]{\scriptsize identité algébrique}};

\node[der, below=1.15cm of born] (lind)
  {Schr\"odinger--\\Lindblad\\[2pt]{\scriptsize décohérence}};

\node[thm, below=1.15cm of lind] (spec)
  {$\rho(T\!\cdot\!D_j) < 1$\\rayon spectral\\[2pt]{\scriptsize Perron--Frobenius}};

\node[thm, below=1.15cm of spec] (t4)
  {T4~: convergence\\$\alpha_k \to 1/2$\\[2pt]{\scriptsize annihilation spectrale}};

\draw[arri] (hilb) -- node[right, font=\scriptsize, text=idgreen]
  {$\sqrt{P}$} (born);
\draw[arrd] (born) -- node[right, font=\scriptsize, text=derorange]
  {$\partial_t\rho$} (lind);
\draw[arrt] (lind) -- node[right, font=\scriptsize, text=thmblue]
  {$\rho < 1$} (spec);
\draw[arrt] (spec) -- node[right, font=\scriptsize, text=thmblue]
  {annihilation} (t4);

% ============================================================
%  ROW BOTTOM: 43 SM observables
% ============================================================
\node[val=5.5cm, below=1.4cm of alpha] (obs)
  {\textbf{43 observables MS}\\$\langle\epsilon\rangle = 0{,}30\%$,
   0~ajusté\\[2pt]{\scriptsize\textsc{val}\enspace(\cite{companion_physics})}};

\draw[arrt] (alpha) -- (obs);
\draw[arrt] (t4.south) |- node[near start, right, font=\scriptsize,
  text=thmblue] {$Q\to 0$} (obs.east);

% ---- Predictions box (red, tabular) ----------------------------
\node[pred=11.0cm, below=1.0cm of obs] (pred)
  {\textbf{\textcolor{predred}{15 prédictions}}
   \enspace{\scriptsize\PREDTAG\enspace(\cite{companion_physics})}\\[4pt]
   {\scriptsize
   \begin{tabular}{@{}r@{\;}l@{\qquad}r@{\;}l@{\qquad}r@{\;}l@{}}
   \multicolumn{2}{@{}l}{\textbf{Niveau 1~: structurelles}} &
   \multicolumn{2}{l}{\textbf{Niveau 2~: numériques}} &
   \multicolumn{2}{l@{}}{\textbf{Niveau 3~: négatives}} \\[2pt]
   P1 & $\nu$ de Dirac       & P5 & $m_{\nu_3} = 0{,}0505$ eV  & P10 & Pas d'axion \\
   P2 & Hiérarchie normale   & P6 & $\sin^2\!\theta_{23}$    & P11 & Pas de SUSY $<100$ TeV \\
   P3 & $\theta_{\rm QCD}=0$ & P7 & $\sin^2\!\theta_{13}$  & P12 & $N_{\rm gen}=3$ exact \\
   P4 & $\delta_{CP}=197^{\circ}$ & P8 & $\lambda_H = 1/8$        & P13 & Pas de désintégration du proton \\
      &                    & P9 & $H_0 = 67{,}41$             & P14 & Pas de WIMP \\
      &                    & P15 & $\alpha_{\rm GW}<10^{-3}$ &     & \\
   \end{tabular}}};
\draw[arrp] (obs) -- (pred);

% ============================================================
%  CROSS-COLUMN LINKS
% ============================================================

% Z/mZ -> Delta_m (residue frequencies)
\draw[cross] (ring.west) -- node[above, font=\scriptsize]
  {$p_r = \#\{n : r_n \equiv r\}/N$} (prob.east);

% CRT -> Hilbert (tensor)
\draw[cross] (crt.east) -- ++(2.0,0) |- node[near end, above,
  font=\scriptsize] {$\mathbb{Z}/m\mathbb{Z}\cong\prod\mathbb{Z}/p\mathbb{Z}
  \;\Leftrightarrow\;\mathcal{H}_m = \bigotimes\mathbb{C}^p$} (hilb.west);

% T_m -> spectral (same operator)
\draw[cross] (tm.east) -- ++(2.5,0) |-
  node[near end, above, font=\scriptsize] {spectre de $\Tm$} (spec.west);

% GFT <-> ring (center)
\draw[cross] (gft.east) -- ++(0.8,0) |- node[near start, above right,
  font=\scriptsize] {} (ring.west);

% holonomy -> Einstein via bridge (identification, proved via OS)
\draw[arrb] (holo.west) -- ++(-2.5,0) |-
  node[near end, above, font=\scriptsize, text=bridgepurple]
  {identification physique (\textsc{thm}, sec.~\ref{sec:bridge})} (einstein.east);

% ============================================================
%  PROOF ROBUSTNESS ANNOTATION
% ============================================================
\begin{scope}[on background layer]
  % Highlight robust nodes (multiple independent proofs)
  \node[draw=thmblue!30, fill=thmblue!3, rounded corners=8pt,
    inner sep=8pt, fit=(shalf)(gn), label={[font=\scriptsize\itshape,
    text=thmblue!60]right:$\geq 3$ preuves indép.}] {};
  \node[draw=idgreen!30, fill=idgreen!3, rounded corners=8pt,
    inner sep=8pt, fit=(fisher)(dkl), label={[font=\scriptsize\itshape,
    text=idgreen!60]left:G1+G3~: 14/14}] {};
  \node[draw=thmblue!30, fill=thmblue!3, rounded corners=8pt,
    inner sep=8pt, fit=(mu15), label={[font=\scriptsize\itshape,
    text=thmblue!60]right:6 preuves indép.}] {};
\end{scope}

% ============================================================
%  LEGEND (compact, bottom of diagram)
% ============================================================
\node[anchor=north, font=\small\bfseries] (leg) at
  ([yshift=-1.2cm]pred.south) {\textsc{Légende}};

\begin{scope}[shift={(leg.south)}, xshift=-4.5cm, yshift=-0.15cm,
  every node/.style={font=\scriptsize, anchor=west}]
  % Left sub-column
  \draw[arrt] (0,0) -- ++(1.0,0) node[right]
    {théorème \THMTAG};
  \draw[arri] (0,-0.4) -- ++(1.0,0) node[right]
    {identité \IDTAG};
  \draw[arrd] (0,-0.8) -- ++(1.0,0) node[right]
    {dérivé \DERTAG};
  \draw[arrp] (0,-1.2) -- ++(1.0,0) node[right]
    {prédiction \PREDTAG};
  % Right sub-column
  \draw[arrb] (5.5,0) -- ++(1.0,0) node[right]
    {pont \BRIDGETAG};
  \draw[cross] (5.5,-0.4) -- ++(1.0,0) node[right]
    {lien inter-colonnes};
  \draw[eqv]  (5.5,-0.8) -- ++(1.0,0) node[right]
    {équivalence ($\equiv$)};
\end{scope}

\end{tikzpicture}%
\end{figure}

\clearpage
\captionof{figure}[Architecture mathématique de la théorie de la persistance]{%
  \textbf{Architecture mathématique de PT.}
  À partir d'un unique théorème ($s = 1/2$), PT engendre trois
  piliers interconnectés~: géométrie de l'information (gauche),
  algèbre du crible (centre) et structure quantique (droite).
  Toutes les flèches portent un statut épistémique explicite. La
  fermeture BA0 (T0, \cite{companion_M3}) élimine tout axiome
  irréductible~; BA5 ($\alphaEM$) est prouvé via la dualité de
  Pontryagin + reconstruction d'OS (\textsc{thm},
  section~\ref{sec:bridge}). Les halos d'arrière-plan marquent les
  n{\oe}uds à plusieurs preuves indépendantes ($\kappa \geq 2$).
  À comparer avec la hiérarchie classique des structures
  mathématiques (groupes, anneaux, corps, espaces métriques,
  espaces de Hilbert)~: PT instancie \emph{simultanément les trois
  branches} à partir d'un unique degré de liberté.}
\label{fig:pt_architecture}

% ============================================================
%  Commentary
% ============================================================
\subsection{Lecture du diagramme}
\label{sec:app_h_reading}

Le diagramme (figure~\ref{fig:pt_architecture}) suit la hiérarchie
classique des structures mathématiques en parallèle, mais en
inverse la logique. Là où le diagramme classique dit «~un espace
de Hilbert \emph{est} un espace de Banach~», notre diagramme dit
«~le crible \emph{engendre} un espace de Hilbert~». Les trois
colonnes correspondent à trois lectures indépendantes du même
objet~$\{\gn\}$~:

\begin{enumerate}[nosep]
\item \textbf{Géométrie de l'information} (gauche)~: les fréquences
  de résidus forment un simplexe de probabilité~$\Delta_m$~;
  l'unicité de Shore--Johnson (G1) sélectionne $\DKL$ comme
  divergence~; l'identité TFS la relie à l'algèbre du crible~;
  l'unicité de \v{C}encov (G3) sélectionne la métrique de Fisher~;
  la métrique de Fisher \emph{est} la limite continue (R50,
  dissoute), donnant les équations d'Einstein (\textsc{thm}, 38/38)
  avec zéro nouveau paramètre.

\item \textbf{Algèbre du crible} (centre)~: la suite des sauts vit
  dans $\mathbb{Z}/m\mathbb{Z}$~; le théorème chinois des restes la
  décompose en corps $\mathbb{Z}/p\mathbb{Z}$~; les axiomes
  C1--C4 forcent le crible d'Ératosthène comme unique solution
  (T6)~; la matrice de transfert~$\Tm$ encode la dynamique~; les
  angles d'holonomie déterminent l'unique point fixe
  $\mustar = 15$ (T5) et donnent $\alphaEM$ via la dualité de
  Pontryagin et la reconstruction d'OS~(BA5, \textsc{thm}).

\item \textbf{Structure quantique} (droite)~: le produit tensoriel
  TCR induit un espace de Hilbert
  $\mathcal{H}_m = \bigotimes \mathbb{C}^p$~; la règle de Born est
  une identité algébrique~; la contraction spectrale ($\rho < 1$)
  assure la convergence~(T4).
\end{enumerate}

Les liens en pointillés inter-colonnes montrent que les trois
colonnes ne sont pas indépendantes~:
$\mathbb{Z}/m\mathbb{Z} \cong \prod\mathbb{Z}/p\mathbb{Z}$ est
simultanément la TCR de la théorie des anneaux (centre) et la
factorisation tensorielle de l'espace de Hilbert (droite)~;
l'identité TFS $\log_2 m = \DKL + H$ relie la géométrie de
l'information (gauche) à l'algèbre du crible (centre).

\subsubsection*{Robustesse des preuves}
\label{sec:app_h_robustness}

Les halos d'arrière-plan du diagramme marquent les n{\oe}uds qui
possèdent au moins deux chemins de preuve disjoints en sommets
depuis les axiomes ($\kappa \geq 2$). Les n{\oe}uds les plus
robustes sont $s = 1/2$ ($\kappa \geq 3$), le couple
Fisher--$\DKL$ (G1+G3, 14/14~tests), et $\mustar = 15$
(6~justifications indépendantes). La fermeture BA0 (T0,
\cite{companion_M3}) élimine tout axiome irréductible~: la théorie
repose sur la seule théorie des nombres. BA5 ($\alphaEM$) porte
$\kappa = 2$ chemins de preuve indépendants (ITPS + Buchstab/OS,
section~\ref{sec:bridge}). L'identité d'holonomie $\sinp{p}$, malgré
sa haute centralité, a $\kappa = 1$~--- une cible stratégique pour
de futures preuves alternatives.


\newpage

\section{Catalogue des transformations PT}
\label{app:transformations}

Cette annexe fournit une référence unifiée pour les dix
transformations qui forment la boîte à outils analytique de la
théorie de la persistance. Chaque transformation capture un aspect
différent de la persistance informationnelle dans le crible~;
ensemble, elles subsument l'analyse de Fourier/Mellin classique
dans le cadre spécifique au crible.

L'unique entrée est $\spt = 1/2$. Aucune transformation n'utilise
de paramètre ajusté.

% ============================================================
\subsection*{Vue d'ensemble}

\begin{center}
\small
\begin{tabular}{clcll}
\toprule
\textbf{\#} & \textbf{Transformation} & \textbf{Dim/K} &
  \textbf{Capte} & \textbf{Outil} \\
\midrule
1 & Persistance $P$ & 2 &
  Secteurs $v_+/v_-$ de $T_3$ & M16 \\
2 & Intercalateur $J$ & --- &
  Échange $v_+ \leftrightarrow v_-$ & M12 \\
3 & Multi-modulaire & 15 &
  Projections mod 3,5,7 (concat.) & M30 \\
4 & Canal CPTP & $3\times 3$ &
  Décohérence quantique & M27 \\
5 & Convolution $\Z/3\Z$ & --- &
  Contraction de Fourier (fusion) & T5 \\
6 & Pontryagin--PT & --- &
  TCR additif $\to$ produit d'Euler & BA5 \\
7 & Tenseur $\mathcal{T}$ & 105 &
  Corrélations inter-modulaires & \textbf{M33} \\
8 & Holonomique $\mathcal{H}$ & $K$ premiers &
  Angles $\sinp{p}$ par premier & \textbf{M34} \\
9 & Décohérence $\mathcal{D}$ & 1 &
  Entropie monotone (théorème H) & \textbf{M35} \\
10 & Levée complexe $\mathcal{W}$ & $\mathbb{C}$ &
  Phase + cohérence sur $\mathcal{C}(s)$ & \textbf{M42} \\
\bottomrule
\end{tabular}
\end{center}

Les descriptions détaillées de chaque transformation, y compris
les équations définissantes, les propriétés, la classification et
les scores de test, sont données dans l'article~IV~\cite{companion_M4}.
On enregistre ici le tableau comparatif et le diagramme de
relations.

% ============================================================
\subsection*{Tableau comparatif}

\begin{center}
\small
\begin{tabular}{lccccc}
\toprule
\textbf{Propriété} & $P$ (M16) & $\mathcal{H}$ (M34) &
  $\mathcal{T}$ (M33) & $\mathcal{D}$ (M35) & Fourier \\
\midrule
Base & $v_\pm$ de $T_3$ & $\sinp{p}$ &
  $\bigotimes_q \mathrm{vec.\,p.}(T_q)$ & entropie & $e^{ikx}$ \\
Paramètre & profondeur $K$ & premiers $p$ & $(K, q_1, q_2, q_3)$ &
  profondeur $K$ & fréq.\ $k$ \\
Dim/K & 2 & $K$ (premiers) & 105 & 1 & $\infty$ \\
Linéaire & OUI & NON & OUI & NON & OUI \\
Multiplicative & NON & approx & NON & NON & OUI \\
Monotone & NON & NON & NON & \textbf{OUI} & NON \\
Inversion & partielle & prod.\ d'Euler & décomp.\ CP & --- & exacte \\
Voit le crible & \textbf{OUI} & \textbf{OUI} & \textbf{OUI} &
  \textbf{OUI} & NON \\
Corr.\ inter-mod. & NON & partielle & \textbf{OUI} & NON & NON \\
Courbure & NON & \textbf{OUI} (Fisher) & NON & NON & NON \\
\bottomrule
\end{tabular}
\end{center}


% ============================================================
\subsection*{Diagramme de relations}

\begin{center}
\begin{tikzpicture}[
  scale=0.70, every node/.append style={transform shape},
  box/.style={draw, thick, rounded corners=3pt, align=center,
    font=\small, inner sep=4pt, minimum height=0.7cm,
    text width=3.2cm, fill=white},
  new/.style={box, draw=thmblue, fill=thmblue!8},
  old/.style={box, draw=black!50, fill=gray!5},
  arr/.style={-{Stealth[length=4pt]}, thick},
]
  % Row 1
  \node[old] (H) at (0,0) {Holonomique $\mathcal{H}$\\(M34)};
  \node[old] (pont) at (-5,0) {Pontryagin\\(BA5)};

  % Row 2
  \node[old] (P) at (-5,-2.5) {Persistance $P$\\(M16)};
  \node[old] (J) at (0,-2.5) {Intercalateur $J$\\(M12)};
  \node[old] (fish) at (4.5,0) {Métrique de Fisher\\(\cite{companion_M2})};

  % Row 3
  \node[old] (MM) at (-5,-5) {Multi-modulaire\\(M30)};
  \node[old] (conv) at (0,-5) {Convolution $\Z/3\Z$\\(T5)};

  % Row 4
  \node[new] (tens) at (-5,-7.5) {Tenseur $\mathcal{T}$\\(\textbf{M33})};
  \node[old] (contr) at (0,-7.5) {Contraction\\de Fourier};

  % Row 5
  \node[new] (D) at (-2.5,-10) {Décohérence $\mathcal{D}$\\(\textbf{M35})};
  \node[new] (thmH) at (-2.5,-12) {Théorème H\\$S_K$ croissant};

  % Arrows
  \draw[arr] (pont) -- (H) node[midway,above,font=\tiny] {dualité};
  \draw[arr] (H) -- (fish) node[midway,above,font=\tiny] {$ds^2$};
  \draw[arr] (H) -- (P) node[midway,left,font=\tiny] {$\sinp{p}$};
  \draw[arr] (P) -- (J) node[midway,above,font=\tiny] {$D_4$};
  \draw[arr] (P) -- (MM);
  \draw[arr] (J) -- (conv);
  \draw[arr] (MM) -- (tens);
  \draw[arr] (conv) -- (contr);
  \draw[arr] (tens) -- (D);
  \draw[arr] (contr) -- (D);
  \draw[arr] (D) -- (thmH);
\end{tikzpicture}
\end{center}


% ============================================================
\subsection*{Pourquoi ces transformations sont originales}

\begin{enumerate}[nosep]
\item Aucune n'est un cas particulier de Fourier ou de Mellin~: la
  base spectrale est celle de $T_q$, non $\{e^{ikx}\}$.
\item Le paramètre est la profondeur du crible~$K$, non une
  fréquence.
\item Elles voient la \emph{structure du crible}~: comment une
  fonction se comporte différemment selon les classes de sauts
  mod~$q$. Fourier/Mellin ne distinguent pas les classes de sauts.
\item Le théorème~H (monotonie de $\mathcal{D}$) est spécifique au
  crible~: pas d'analogue dans la théorie de Fourier classique.
\item Le rang tensoriel de $\mathcal{T}$ mesure une quantité absente
  de l'analyse harmonique classique~: la complexité des corrélations
  inter-modulaires.
\item La métrique de Fisher de $\mathcal{H}$ émerge naturellement
  des angles d'holonomie, reliant la géométrie différentielle à la
  théorie des nombres.
\end{enumerate}


% ============================================================
\subsection*{Scores de tests}

\begin{center}
\begin{tabular}{lc}
\toprule
\textbf{Outil} & \textbf{PASS / TOTAL} \\
\midrule
M33 (Tenseur $\mathcal{T}$) & 10/10 \\
M34 (Holonomique $\mathcal{H}$) & 13/13 \\
M35 (Décohérence $\mathcal{D}$) & 11/11 \\
M42 (Levée complexe $\mathcal{W}$) & 14/14 \\
\midrule
\textbf{Total nouvelles transformations} & \textbf{48/48} \\
\bottomrule
\end{tabular}
\end{center}


\newpage

\section{Limite inductive et reconstruction d'OS}
\label{app:os_limit}

Cette annexe fournit la preuve complète de la construction de la
limite inductive qui produit une TQC wightmanienne continue à
partir des opérateurs de transfert finis $T_m$ du crible.

% ============================================================
\subsection{Mise en place~: la TQC au niveau fini}
\label{sec:os_setup}

À chaque module sans facteur carré $m = p_1 \cdots p_k$, le crible
définit~:
\begin{itemize}[nosep]
\item Un espace de Hilbert $H_m = \mathbb{C}^{\phi(m)}$
  (dimension = indicateur d'Euler de $m$).
\item Un opérateur de transfert $T_m$ (stochastique-en-ligne,
  matrice $\phi(m) \times \phi(m)$).
\item Un état du vide $|\Omega_m\rangle = \sqrt{\pi_m}$
  (distribution stationnaire).
\item Des fonctions de Schwinger
  $S_n^{(m)}(k_1, \ldots, k_n) =
  \langle\Omega_m| \hat\phi\, T_m^{k_1} \hat\phi\, T_m^{k_2}
  \cdots T_m^{k_n} \hat\phi |\Omega_m\rangle$.
\end{itemize}

\subsection{Axiomes d'OS au niveau fini}
\label{sec:os_finite}

\begin{theoreme}[OS3~: positivité par réflexion~--- uniforme]
\label{thm:app_OS3}
Pour tout $m$ sans facteur carré et tout premier $p \geq 3$~:
\begin{enumerate}[nosep]
\item $M_p = T_p^{\mathsf{T}} T_p$ est une matrice de Gram
  ($M_p = X^{\mathsf{T}} X$ avec $X = T_p$, donc $M_p \geq 0$).
\item $M_m = \bigotimes_{p | m} M_p$ est positive semi-définie
  (produit de Kronecker de matrices SDP est SDP).
\item La fonction de Schwinger à deux points satisfait
  $S_2^{(m)}(k) = \langle\Omega_m| \hat\phi\, T_m^k \hat\phi
  |\Omega_m\rangle \geq 0$ pour tout $k \geq 0$.
\end{enumerate}
\begin{proof}
(1)~Pour toute matrice réelle $X$, $X^{\mathsf{T}} X$ est SDP~:
$v^{\mathsf{T}} X^{\mathsf{T}} X v = \|Xv\|^2 \geq 0$. Appliquer
avec $X = T_p$.

(2)~Si $A \geq 0$ et $B \geq 0$ (SDP), alors $A \otimes B \geq 0$~:
pour tout vecteur $w$, $w^{\mathsf{T}} (A \otimes B) w =
\sum_{ij} A_{ij} (w_i^{\mathsf{T}} B\, w_j) \geq 0$ par le
théorème du produit de Schur. La récurrence sur le nombre de
facteurs donne $M_m = \bigotimes_p M_p \geq 0$.

(3)~$S_2^{(m)}(k) = \pi_m^{\mathsf{T}} \mathrm{diag}(\hat\phi)\,
T_m^k\, \mathrm{diag}(\hat\phi)\, \pi_m$. Comme $T_m^k$ est une
matrice stochastique à entrées non négatives et $\pi_m, \hat\phi$
ont des composantes non négatives, le produit est non négatif.
\end{proof}
\end{theoreme}

Les axiomes d'OS restants (OS1~: invariance euclidienne, OS2~:
symétrie, OS4~: propriété d'agrégat) découlent de~:
\begin{itemize}[nosep]
\item OS1~: la structure produit TCR assure l'invariance par
  translation dans chaque direction $\mathbb{Z}/p\mathbb{Z}$.
\item OS2~: les fonctions de Schwinger sont symétriques sous
  permutation des arguments temporels (la matrice de transfert
  commute avec l'action du groupe cyclique).
\item OS4~: le gap spectral $\gamma_K > 0$ (\cite{companion_M3})
  implique le regroupement exponentiel~:
  $|S_2^{(m)}(k) - \langle\hat\phi\rangle^2| \leq C\,|\lambda_2|^k
  \to 0$.
\end{itemize}

\subsection{La limite inductive}
\label{sec:os_inductive}

\begin{theoreme}[Limite inductive via contraction de Buchstab]
\label{thm:app_inductive}
La famille $\{(H_m, T_m, \Omega_m)\}_{m \text{ sans fact.\ carré}}$
admet une limite inductive $(H_\infty, T_\infty, \Omega_\infty)$
au sens des espaces de Hilbert projectifs.

\begin{proof}
La preuve procède en trois étapes.

\textbf{Étape 1~: système projectif.}
Pour $m_1 | m_2$ (c'est-à-dire $m_1$ divise $m_2$), la projection
TCR
$\pi_{m_2 \to m_1} : \mathbb{Z}/m_2\mathbb{Z} \to
\mathbb{Z}/m_1\mathbb{Z}$ induit une application linéaire bornée
$\Phi_{m_2 \to m_1} : H_{m_2} \to H_{m_1}$ définie par
$\Phi(f)(r) = \sum_{r' : r' \equiv r \pmod{m_1}} f(r')$. Ces
applications satisfont la condition de compatibilité
$\Phi_{m_2 \to m_1} \circ \Phi_{m_3 \to m_2} =
\Phi_{m_3 \to m_1}$ pour $m_1 | m_2 | m_3$ (fonctorialité de la
projection TCR). La famille $(H_m, \Phi_{m_2 \to m_1})$ est un
système projectif d'espaces de Hilbert.

\textbf{Étape 2~: contraction de Buchstab.}
La fonction de Buchstab $\omega(u)$ satisfait
$\omega(u) \to e^{-\gamma}$ lorsque $u \to \infty$ (résultat
classique, Buchstab 1937). Au niveau de crible $K$ avec primorielle
$P_K = \prod_{p \leq p_K} p$, la densité de survie est
$\rho_K = \prod_{p \leq p_K}(1 - 1/p) \sim e^{-\gamma}/\ln P_K
\to 0$. La norme de la projection
$\|\Phi_{P_{K+1} \to P_K}\|$ satisfait
$\|\Phi\| \leq \sqrt{(p_{K+1}-1)/p_{K+1}} < 1$ (contraction).
Donc les fonctions de Schwinger $S_n^{(P_K)}$ convergent lorsque
$K \to \infty$~:
$\|S_n^{(P_{K+1})} - S_n^{(P_K)}\| \leq C_n \cdot \rho_K^{1/2} \to 0$.

\textbf{Étape 3~: reconstruction d'OS.}
Par le théorème~\ref{thm:app_OS3}, OS3 (positivité par réflexion)
est vrai à chaque niveau fini. Par l'étape~2, les fonctions de
Schwinger convergent. Par le théorème de reconstruction
d'Osterwalder--Schrader~\cite{OsterwalderSchrader1973,OsterwalderSchrader1975}~:
une suite de fonctions de Schwinger satisfaisant OS1--OS4 et
convergeant au sens des distributions produit une TQC wightmanienne
$(H_\infty, \Omega_\infty, W_n)$ satisfaisant les axiomes
wightmaniens (covariance de Lorentz, condition spectrale,
localité, complétude).

L'espace de Hilbert limite inductive est
$H_\infty = \varprojlim H_{P_K}$ (limite projective dans la
catégorie des espaces de Hilbert), avec
$T_\infty = \varprojlim T_{P_K}$ et
$\Omega_\infty = \varprojlim \Omega_{P_K}$.
\end{proof}
\end{theoreme}

\begin{remarque}[Vérification ITPS]
\label{rem:ITPS}
Une vérification indépendante utilise le produit tensoriel infini
d'espaces de Hilbert (ITPS) de von Neumann~:
$H_\infty = \bigotimes_{p \geq 3}^{\mathrm{vN}} H_p$. La
construction ITPS n'exige pas une dimension fixée des espaces
composantes (Guichardet 1966~\cite{Guichardet1966},
Araki--Woods 1968~\cite{ArakiWoods1968})~: elle s'applique à
$\dim H_p = p - 1$ (croissant avec $p$). Les vecteurs
stabilisants sont les distributions stationnaires
$\Omega_p = \sqrt{\pi_p}$ à chaque premier. L'ITPS et la limite
projective de Buchstab produisent le même $H_\infty$ (toutes deux
sont des complétions de la limite inductive algébrique par
rapport au produit scalaire hérité des fonctions de Schwinger).

Cela fournit un chemin indépendant vers la limite continue,
contournant entièrement le théorème de reconstruction d'OS. Les
deux voies (OS + Buchstab vs.\ ITPS + Guichardet) ne partagent
aucune étape commune au-delà des données de niveau fini.
\end{remarque}

\begin{remarque}[Subtilité restante]
\label{rem:os_subtlety}
La contraction de Buchstab (étape~2) garantit la convergence des
fonctions de Schwinger en norme, ce qui est plus fort que la
convergence au sens des distributions. Toutefois, le taux de
convergence dépend du gap spectral $\gamma_K$, qui est contrôlé
par T4 (\cite{companion_M3}). La preuve est donc logiquement
dépendante de T4 (via la borne sur le gap spectral), mais d'aucune
hypothèse physique.
\end{remarque}

\bibliographystyle{plainnat}
\bibliography{references}

\end{document}
